\title{Plantilla TFM: Euclides}

% ======================================================================
% PLANTILLA LaTeX — Euclides
% ----------------------------------------------------------------------
% Motor requerido: LuaLaTeX o XeLaTeX (se usa fontspec para la fuente).
%   Compilar con:  lualatex nombre.tex   (dos veces, por las referencias
%                  cruzadas de ecuaciones, definiciones, teoremas, etc.)
% ======================================================================
\documentclass[11pt,a4paper,twoside,openright]{book}    
\usepackage{amsmath,amsfonts,amssymb,amsthm,amscd} 


% se copian las variables más relevantes
\makeatletter
\let\newtitle\@title
\let\newauthor\@author
\let\newdate\@date
\makeatother


% ----------------------------------------------------------------------
% Fuente y tamaño de letra
% ----------------------------------------------------------------------
\usepackage{fontspec}
\usepackage{newcomputermodern} % Tipo de letra, New Computer Modern

%
%% Tipo de letra de las letras capitales.
%

%\newfontfamily\chapnumfont{arenq.regular}[Path = ./, Extension = .otf] % NO
%\newfontfamily\chapnumfont{manbow.lines-regular}[Path = ./, Extension = .otf] % de momento, la que más me gusta
%\newfontfamily\chapnumfont{evensong}[Path = ./, Extension = .otf]
%\newfontfamily\chapnumfont{berylium.bold}[Path = ./, Extension = .otf]
%\newfontfamily\chapnumfont{library-3-am.3am}[Path = ./, Extension = .otf] % este 
%\newfontfamily\chapnumfont{foglihtenno04.regular}[Path = ./, Extension = .otf]
%\newfontfamily\chapnumfont{labyrinth.regular}[Path = ./, Extension = .otf]
%\newfontfamily\chapnumfont{baflion-sans.black}[Path = ./, Extension = .otf]
%\newfontfamily\chapnumfont{atavian.regular}[Path = ./, Extension = .otf]
%\newfontfamily\chapnumfont{stripescaps.regular}[Path = ./, Extension = .ttf] % non me gusta e non ten tilde
%\newfontfamily\chapnumfont{dita-sweet.regular}[Path = ./, Extension = .otf]
%\newfontfamily\chapnumfont{ZnikomitNo25}[Path = ./, Extension = .otf]
\newfontfamily\chapnumfont{Archistico_Bold}[Path = ./, Extension = .ttf]
%\newfontfamily\chapnumfont{Roke1984}[Path = ./, Extension = .otf]
%\newfontfamily\chapnumfont{SketchiquaB}[Path = ./, Extension = .ttf]
%\newfontfamily\chapnumfont{WelcomeToFuture1-E4Xnl}[Path = ./, Extension = .ttf]
%\newfontfamily\chapnumfont{CrazyConstructor}[Path = ./, Extension = .ttf]
%\newfontfamily\chapnumfont{Limelight-Regular}[Path = ./, Extension = .ttf]



% interlineado
\renewcommand{\baselinestretch}{1.25}
\setlength{\parskip}{1mm}
\raggedbottom\setlength{\baselineskip}{1mm}
\allowdisplaybreaks



% ----------------------------------------------------------------------
% Idioma
% ----------------------------------------------------------------------
\usepackage[spanish,es-lcroman,es-noquoting]{babel} % Español



% ----------------------------------------------------------------------
% Color primario
% ----------------------------------------------------------------------
\usepackage{xcolor}
\definecolor{colorprimario}{HTML}{FF4000}%{FF461A}



% ----------------------------------------------------------------------
% Geometría de la Página
%-----------------------------------------------------------------------
\usepackage[papersize={210mm,297mm}, left=2.5cm, right=2.5cm, top=3.5cm, bottom=3.5cm]{geometry}


% ----------------------------------------------------------------------
% Paquetes auxiliares
% ----------------------------------------------------------------------
\usepackage{array}       % para el separador "!{...}" en tablas
\usepackage{tabularx}     % columna elástica "X" 
\usepackage{framed}       % entorno "leftbar": regla vertical izquierda
\usepackage{graphicx}
\usepackage{caption}
\usepackage{tikz}
\usepackage{calc}
\usepackage{ulem} % para poder subrayar correctamente
\usepackage{scrextend} % para comandos como addmargin
\usepackage{needspace} % para evitar roturas de cajas no bonitas



% ----------------------------------------------------------------------
% Apéndices
% ----------------------------------------------------------------------
%% Formato: 
%% - Los apéndices se disponen en una Parte separada, 
%%
%% - con una página indicando el inicio con `\appendixtitle` 
%%   escrito con la primera letra en mayúscula y de `colorprimario`. 
%%
%% - El título está subrayado y centrado en el
%%   centro de la página de estilo `empty`, en una página a la izquierda.
%%
%% - Siguiente página en blanco.
%

\usepackage[toc,page]{appendix}

%
%% comandos de formateo
%

\newcommand*{\appendixfontsize}{50pt} % tamaño del título de la sección de anexos
\newcommand*{\appendixtitle}{Apéndices} % título de la sección de Anexos


% TRABAJO INTERNO: -----------------------------------------------------
\newlength{\appendixfontsizeInternal}

%
%% Formato del título de la zona de apéndices
%
\makeatletter
% 1. El motor interno: toma la primera letra (#1), el resto (#2), y se detiene en \@nil
\def\@appendixformat#1#2\@nil{%
	% Primera letra (#1) con su estilo y subrayado propio
	\uline{{\fontsize{1.8\appendixfontsizeInternal}{1.8\appendixfontsizeInternal}\chapnumfont\selectfont\bfseries\color{colorprimario}#1}%
	% Resto del texto (#2) con su propio estilo
	{\fontsize{\appendixfontsize}{\appendixfontsize}\selectfont\bfseries#2}}%
}

% 2. El comando público: expande tu variable antes de enviarla al motor interno
\newcommand{\formattingappendixname}[1]{%
	\expandafter\@appendixformat#1\@nil
}

% 
%% Definición del apéndice como tal
%
\gdef\appendix{%
	\setlength{\appendixfontsizeInternal}{\appendixfontsize}
	\renewcommand{\ULthickness}{1.5pt}
	\cleardoublepage
	\thispagestyle{empty} % Quita encabezados y pies de página en esta portada
	
	% Página con "Apéndices" alineado a la izquierda y centrado verticalmente
	\vspace*{\fill}
	\begin{center}
		\formattingappendixname{\appendixtitle} % formateo del título
	\end{center}
	\vspace*{\fill}
	
	\clearpage
	\thispagestyle{empty}
	\null % Inserta una página en blanco obligatoria
	\clearpage
	
	% Restablece los contadores para los capítulos del apéndice
	\setcounter{chapter}{0}%
	\setcounter{section}{0}%
	\gdef\@chapapp{\appendixname}%
	\gdef\thechapter{\@Alph\c@chapter}%
}
\makeatletter



% ----------------------------------------------------------------------
% Índice de contenidos
% ----------------------------------------------------------------------
\usepackage{titletoc}

%
%% comandos de formato
%

\renewcommand\contentsname{Índice} % nombre del título

\newcommand{\contentstitlesize}{50pt} % tamaño de letra del título




% TRABAJO INTERNO ------------------------------------------------------
\newlength{\contentstitlesizeInternal}


\titlecontents{chapter}[0pt]
{\addvspace{1em}\bfseries}                 % antes de la entrada
{\textcolor{colorprimario}\thecontentslabel\quad}        % con número
{}                                            % sin número (no se usa aquí)
{\leavevmode                                  % <- después del título:
	\hfill\kern0pt    %    
	\textcolor{colorprimario}{\thecontentspage}} %    página en naranja


\titlecontents{section}[1.5em]
{\addvspace{0.10em}}
{\textcolor{colorprimario}\thecontentslabel.\quad}
{}
{\leavevmode
	~\leaders\hrule height 0.5pt\hfill\kern0pt ~   % línea recta de relleno
	\textcolor{colorprimario}{\thecontentspage}}


\titlecontents{subsection}[3em]
{\addvspace{0.1em}\small}
{\thecontentslabel.\quad}
{}
{\leavevmode
	~\leaders\hrule height 0.5pt\hfill\kern0pt ~
	\textcolor{colorprimario}{\thecontentspage}}




\let\oldtableofcontents\tableofcontents % guardamos la anterior forma de table of contents.

% Formateo del título
\makeatletter
\renewcommand{\tableofcontents}{%
	\begingroup
	\hypersetup{linkcolor=black}% cambiar, de forma temporal, el color de los títulos.
	\renewcommand{\@makeschapterhead}[1]{%
		\setlength{\contentstitlesizeInternal}{\contentstitlesize}
		\renewcommand{\arrayrulewidth}{1.5pt}
		\vspace*{-5mm}%
		\noindent
		\begin{tabularx}{\textwidth}{ l @{\vline \phantom{M}} X}
			 & \vphantom{{\fontsize{1.8\contentstitlesizeInternal}{1.8\contentstitlesizeInternal}\selectfont M}} \textbf{{\fontsize{1.2\contentstitlesizeInternal}{1.2\contentstitlesizeInternal}\chapnumfont\selectfont\color{colorprimario}Í}\scshape\fontsize{\contentstitlesizeInternal}{\contentstitlesizeInternal}\selectfont ndice} \\
			\llap{\raisebox{1ex}{\rule{20cm}{\arrayrulewidth}}\hspace{-\textwidth}} &\\ &
		\end{tabularx}
		\renewcommand{\arrayrulewidth}{1pt}
	}%
	\oldtableofcontents
	\endgroup
}
\makeatother



% ----------------------------------------------------------------------
% Cabecera: 
% ----------------------------------------------------------------------
\usepackage{fancyhdr}
\pagestyle{fancy}
\fancyhf{}
\renewcommand{\headrulewidth}{0.5pt}
\renewcommand{\footrulewidth}{0pt}

\renewcommand{\chaptermark}[1]{\markboth{\textbf{{\color{colorprimario}\thechapter} #1}}{}} % formato del título de capítulo

\fancypagestyle{fancy}{	
	\fancyhf{}
	\renewcommand{\headrulewidth}{0.5pt}
	\renewcommand{\footrulewidth}{0pt}
	\fancyhead[RO]{\bfseries\fontsize{11pt}{13pt}\selectfont \newtitle}
	\fancyhead[LE]{\fontsize{11pt}{13pt}\selectfont \leftmark}
	\fancyhead[RE,LO]{\color{colorprimario}\bfseries\thepage}
}

\fancypagestyle{plain}{%
	\fancyhf{}
	\renewcommand{\headrulewidth}{0pt}
%	\renewcommand{\footrulewidth}{0pt}
	\fancyfoot[C]{\color{colorprimario}\bfseries\thepage}
}


% ----------------------------------------------------------------------
% Títulos de capítulo y secciones
% ----------------------------------------------------------------------
\usepackage{titlesec}


% titulos de secciones
\titleformat{\section}
{\normalfont\Large\bfseries}
{\color{colorprimario}\thesection}
{0.5em}
{\hyphenpenalty=10000\exhyphenpenalty=10000\relax}

\titleformat{\subsection}
{\normalfont\large\bfseries}
{\color{colorprimario}\thesection.\arabic{subsection}}
{0.5em}
{\hyphenpenalty=10000\exhyphenpenalty=10000\relax}


% títulos de capítulo
\makeatletter

%
%% Comandos de Formato
%

\newcommand*{\chaptertitlesize}{25pt} % definición del tamaño de letra en los títulos
\newcommand*{\chapternumsize}{90pt} % definición del tamaño de fuenta del número de capítulo

\newcommand{\chaptersubsep}{10mm}

% TRABAJO INTERNO ------------------------------------------------------
\newlength{\chaptertitlesizeInternal}
\setlength{\chaptertitlesizeInternal}{\chaptertitlesize}

\newlength{\chapternumsizeInternal}
\setlength{\chapternumsizeInternal}{\chapternumsize}

\newlength{\chaptersubsepInternal}
\setlength{\chaptersubsepInternal}{\chaptersubsep}

\renewcommand{\tabularxcolumn}[1]{b{#1}}

\renewcommand{\@makechapterhead}[1]{%
	\setlength{\chaptertitlesizeInternal}{\chaptertitlesize}
	\setlength{\chapternumsizeInternal}{\chapternumsize}
	\setlength{\chaptersubsepInternal}{\chaptersubsep}
	\renewcommand{\arrayrulewidth}{1.5pt}
	\noindent
	\begin{tabularx}{\textwidth}{ l @{\vline \phantom{M}} X}
%		\llap{\raisebox{1ex}{\rule{25cm}{\arrayrulewidth}}\hspace{-0pt-\widthof{\fontsize{90}{90}\chapnumfont\selectfont\thechapter}}\hspace{-3mm}} & \\ % línea superior en el número
		{\fontsize{\chapternumsize}{\chapternumsize}\chapnumfont\selectfont \color{colorprimario} \thechapter}
		 \vphantom{{\fontsize{1.3\chapternumsizeInternal}{1.3\chapternumsizeInternal}\selectfont \thechapter}} 
		 & \textbf{\scshape\fontsize{\chaptertitlesize}{\chaptertitlesize}\selectfont #1} \\
		\llap{\raisebox{1ex}{\rule{20cm}{\arrayrulewidth}}\hspace{-\textwidth}} & \\ &
	\end{tabularx}
	\vspace{\chaptersubsepInternal}
	\renewcommand{\arrayrulewidth}{1pt}
}

% título sin numerar
\renewcommand{\@makeschapterhead}[1]{%
	\setlength{\chaptertitlesizeInternal}{\chaptertitlesize}
	\setlength{\chapternumsizeInternal}{\chapternumsize}
	\setlength{\chaptersubsepInternal}{\chaptersubsep}
	\renewcommand{\arrayrulewidth}{1.5pt}
	\par\noindent
	\begin{tabularx}{\textwidth}{ l @{\vline \phantom{M}} X}
	%	{\fontsize{90}{90}\selectfont \color{colorprimario} \thechapter} 
	%\llap{\raisebox{1ex}{\rule{25cm}{\arrayrulewidth}}\hspace{-3mm}} & \\
		 \vphantom{{\fontsize{\dimexpr 25pt+\chaptertitlesizeInternal\relax}{\dimexpr 25pt+\chaptertitlesizeInternal\relax}\selectfont M}}
		  & \textbf{\scshape\fontsize{\chaptertitlesizeInternal}{\chaptertitlesizeInternal}\selectfont #1} \\
		\llap{\raisebox{1ex}{\rule{20cm}{\arrayrulewidth}}\hspace{-\textwidth}} &\\ &
	\end{tabularx}
	\vspace{\chaptersubsepInternal}
	\renewcommand{\arrayrulewidth}{1pt}
}
\makeatother


%\titlespacing*{\chapter}{0pt}{-1.5cm}{2cm}
\titlespacing*{\section}{0pt}{0.25cm}{0.25cm}



% ----------------------------------------------------------------------
% Notas al pie coloreadas
% ----------------------------------------------------------------------
\renewcommand{\thefootnote}{\textcolor{colorprimario}{\bfseries\arabic{footnote}}}


% ----------------------------------------------------------------------
% Enlaces coloreados y en negrita
% ----------------------------------------------------------------------
\usepackage[colorlinks=true, linkcolor=colorprimario, citecolor=colorprimario,
urlcolor=colorprimario, filecolor=colorprimario,backref=page]{hyperref}
\let\oldhref\href
\renewcommand{\href}[2]{\oldhref{#1}{\textbf{#2}}}



% ----------------------------------------------------------------------
% Figuras y tablas
% ----------------------------------------------------------------------
\DeclareCaptionLabelFormat{coloreado}{\textcolor{colorprimario}{\textbf{#1~#2}}}
\captionsetup{
	labelformat=coloreado, 
	labelsep=colon}



% ----------------------------------------------------------------------
% Ecuaciones: numerar solo las que están referenciadas + número coloreado
% ----------------------------------------------------------------------
% \usepackage{mathtools}


% IMPORTANTE: el tagform coloreado se define ANTES de activar
% showonlyrefs. mathtools guarda una copia de \tagform@ en el momento
% de activar showonlyrefs y la envuelve con su lógica de "solo
% mostrar si está referenciada"; si se redefine \tagform@ DESPUÉS, se
% destruye esa lógica en vez de envolverla (y el número sale mal).
\makeatletter
\renewcommand{\tagform@}[1]{\maketag@@@{\color{colorprimario}\bfseries(#1)}}
\makeatother



% ----------------------------------------------------------------------
% Bibliografía
% ----------------------------------------------------------------------
% Cambia el separador cuando solo hay dos páginas (ej: 1 y 2)
\renewcommand*{\backreftwosep}{ y }

% Cambia el separador final cuando hay tres o más páginas (ej: 1, 2 y 3)
\renewcommand*{\backreflastsep}{ y }

% Desactiva el formato por defecto
\renewcommand*{\backref}[1]{}

% Define el nuevo formato
%\renewcommand*{\backrefalt}[4]{%
%	\ifcase #1 %
%	% Caso 0: No se ha citado en el texto (puedes dejarlo vacío)
%	\or
%	% Caso 1: Citado en exactamente una página
%	\\ \footnotesize \raggedleft  \textit{#2.}%
%	\else
%	% Caso 2: Citado en múltiples páginas
%	\par \hfill \footnotesize  \texttt{pp. #2.}%
%%	\begin{flushright}
%%		\footnotesize {\textit{Citado en las páginas #2.}}%
%%	\end{flushright} 
%	\fi
%}

% -----
% Titulo: backreferences a las páginas donde las referencias bibliográficas aparecen. 
% Formato:  recuadro grisáceo con las páginas en `colorprimario`
% -----
\makeatletter
\renewcommand*{\backrefalt}[4]{%
	\ifnum#1=0\else
	\nobreak\hfill\break
	\vspace{-25pt}%
	\begin{flushright}
		\def\@lastpage{}%
		\@for\@page:=#2\do{%
			\ifx\@page\@lastpage\else
			\begingroup
			\hypersetup{linkcolor=colorprimario}%
			\setlength{\fboxsep}{2pt}%
			\colorlet{backrefbg}{gray!15}
			\colorlet{backrefborder}{gray!55}
			\fcolorbox{backrefborder}{backrefbg}
			{\footnotesize\bfseries\@page}%
			\endgroup
			\hspace{3pt}%
			\let\@lastpage\@page
			\fi
		}%
	\end{flushright}%
	\fi
}
\makeatother


% cuadro al rededor de las referencias bibliográficas
%\newcommand{\boxcite}[1]{%
%	\begingroup
%	\setlength{\fboxsep}{2.5pt}%
%	\setlength{\fboxrule}{0.8pt}%
%	\fcolorbox{black!20}{white}{%
%		\textcolor{black}{\footnotesize\bfseries#1}%
%	}%
%	\endgroup
%}
%
%\makeatletter
%\renewcommand{\@cite}[2]{%
%	\boxcite{[#1\if@tempswa , #2\fi]}%
%}
%\makeatother

\emergencystretch=5em




% ----------------------------------------------------------------------
% Entornos: Definición, Teorema, Proposición, Lema
% ----------------------------------------------------------------------


% Regla vertical izquierda reutilizable (definiciones/teoremas/.../demos)
\renewenvironment{leftbar}{%
	\def\FrameCommand{{\color{black}\vrule width 0.5pt}\hspace{3mm}}%
	\vspace{-9mm}
	\MakeFramed{\advance\hsize-\width \FrameRestore}%
	\vphantom{M}
	\vspace{1mm}\\
}{\endMakeFramed}



\newlength{\anchocaja}
\newcommand{\openEuclidOverlayleft}[2]{%
	\par\addvspace{1.1em}
	\renewcommand{\arrayrulewidth}{0.5pt}
	\renewcommand{\arraystretch}{0.5}
	\setlength{\anchocaja}{\widthof{MM #1 | #2 MM}}
	\ifdim\anchocaja>\textwidth
	\setlength{\anchocaja}{\textwidth}
	\fi
	\setbox0=\hbox{%
		\begin{tabularx}{\anchocaja}{r|X}
			\vphantom{\Large M}~ \color{colorprimario}\bfseries #1 & #2 \\[0ex]
			\hline
			&
		\end{tabularx}%
	}%
	\Needspace{\dimexpr\ht0+\dp0+3\baselineskip\relax}
	\hspace{-2em}\box0
	\par\nobreak
%	\vskip -9mm\nobreak
	\begin{leftbar}
	}
	\newcommand{\closeEuclidOverlayleft}{%
	\end{leftbar}
	\renewcommand{\arraystretch}{1}
	\addvspace{1.1em}
}



% para las notas

\newlength{\anchocajanota}
\newcommand{\openEuclidOverlayright}[2]{%
	\par\addvspace{1.1em}
	\renewcommand{\arrayrulewidth}{0.5pt}
	\renewcommand{\arraystretch}{0.5}
	\setlength{\anchocajanota}{\textwidth}
%	\hspace{-2em}
	\setbox0=\hbox{
		\begin{tabularx}{\anchocajanota}{X|l}
			\flushright \color{colorprimario}\bfseries #1 & #2 \\[0ex]
			\hline
			&
		\end{tabularx}
	}
	\Needspace{\dimexpr\ht0+\dp0+2\baselineskip\relax}
	\hspace{-2.2em}\box0
	\par\nobreak
%	\nobreak\hfill
	\begin{leftbar}
	}
	\newcommand{\closeEuclidOverlayright}{%
	\end{leftbar}
	\renewcommand{\arraystretch}{1}
	\addvspace{1.1em}
}




% definición de los entornos
\usepackage{etoolbox} % para \ifstrequal

\newcommand{\newenv}[4]{%
	% #1 = identificador, #2 = ámbito, #3 = título, #4 = left|right
	\ifstrequal{#2}{chapter}{%
		\newcounter{#1}[chapter]%
		\expandafter\renewcommand\csname the#1\endcsname{\thechapter.\arabic{#1}}%
	}{%
		\ifstrequal{#2}{section}{%
			\newcounter{#1}[section]%
			\expandafter\renewcommand\csname the#1\endcsname{\thechapter.\arabic{section}.\arabic{#1}}%
		}{%
			\ifstrequal{#2}{subsection}{%
				\newcounter{#1}[subsection]%
				\expandafter\renewcommand\csname the#1\endcsname{\thechapter.\arabic{section}.\arabic{subsection}.\arabic{#1}}%
			}{%
				\newcounter{#1}%
				\expandafter\renewcommand\csname the#1\endcsname{\arabic{#1}}%
	}}}%
	\newenvironment{#1}[1][]{%
		\refstepcounter{#1}%
		\csname openEuclidOverlay#4\endcsname{#3~\csname the#1\endcsname}{##1}%
	}{\csname closeEuclidOverlay#4\endcsname}%
}



% ----------------------------------------------------------------------
% Demostración
% ----------------------------------------------------------------------
\renewcommand{\proofname}{Demostración}
\renewenvironment{proof}[1][\proofname]{%
	\addvspace{1em}\par
	\Needspace{3\baselineskip}
	\pushQED{\qed}%
	\noindent{\color{colorprimario}\bfseries #1}%
	\begin{addmargin}{3mm}
	}{%
		\popQED
	\end{addmargin}
	\addvspace{1em}
}


% ----------------------------------------------------------------------
% Blóques de código
% ----------------------------------------------------------------------

%\usepackage{fancyvrb}
%\fvset{
%	fontsize=\small,
%	frame=single,
%	framesep=3mm,
%	rulecolor=\color{black},  % color del marco
%	numbers=left,                      % numeración de líneas
%	numbersep=5pt,
%	formatcom=\color{black},           % color del texto en sí
%}

\usepackage{listings}

\lstdefinestyle{miR}{
	language=R,
	basicstyle=\ttfamily\small,
	keywordstyle=\color{colorprimario}\bfseries,   
	commentstyle=\color{gray}\itshape,  % # comentarios
	stringstyle=\color{blue}, % "cadenas"
	numbers=left, numbersep=6pt, numberstyle=\tiny\color{gray},
	showstringspaces=false, breaklines=true, xleftmargin=-1.2em,
}


%\usepackage{tcolorbox}
%\tcbuselibrary{listings, breakable}
%
%\newtcblisting{codigo}{
%	breakable,
%	colback=white, 
%	colframe=black,
%	boxrule=0.6pt, 
%	listing only,
%	listing options={style=miR, basicstyle=\ttfamily\small, breaklines=true},
%%	left=3mm, right=3mm, top=1mm, bottom=1mm,
%}


\usepackage{tikz}
\usepackage[many]{tcolorbox}
\tcbuselibrary{listings,skins, breakable}

% Definimos el nuevo entorno para la caja
\newtcblisting{codigo}{%
	breakable,
	enhanced,
	frame hidden,       % Ocultamos el marco normal
	colback=white,      % Color de fondo de la caja
	top=1mm, bottom=1mm, left=4mm, right=4mm, % Márgenes internos
	listing only,
	listing options={style=miR, basicstyle=\ttfamily\small, breaklines=true},
	overlay={
		\def\ext{2mm}     % Define cuántos milímetros sobresalen las líneas
		
		% Línea superior
		\draw[thick, black] ([xshift=-\ext]frame.north west) -- ([xshift=\ext]frame.north east);
		% Línea inferior
		\draw[thick, black] ([xshift=-\ext]frame.south west) -- ([xshift=\ext]frame.south east);
		% Línea izquierda
		\draw[thick, black] ([yshift=\ext]frame.north west) -- ([yshift=-\ext]frame.south west);
		% Línea derecha
		\draw[thick, black] ([yshift=\ext]frame.north east) -- ([yshift=-\ext]frame.south east);
	},
}



% ----------------------------------------------------------------------
% glosario
% ----------------------------------------------------------------------
\usepackage[toc]{glossaries}
\setglossarystyle{listhypergroup}









% ejemplo de un enviroment Definición, Teorema, Proposición, Lema
\newenv{definition}{chapter}{Definición}{left} 
\newenv{theorem}{chapter}{Teorema}{left}
\newenv{proposition}{chapter}{Proposición}{left}
\newenv{lemma}{chapter}{Lema}{left}
\newenv{nota}{chapter}{Nota}{right}




\newcommand{\vertiii}[1]{{\left\vert\kern-0.25ex\left\vert\kern-0.25ex\left\vert #1 
		\right\vert\kern-0.25ex\right\vert\kern-0.25ex\right\vert}}
\def\N{\ensuremath{\mathbb{N}}} %-- la N de los naturales
\def\R{\ensuremath{\mathbb{R}}} %-- La R de los reales.
\def\C{\ensuremath{\mathbb{C}}} %-- La C de los complejos.
\def\Q{\ensuremath{\mathbb{Q}}}
\def\K{\ensuremath{\mathbb{K}}}

\loadglsentries{glosario-prueba.tex}
\makeglossaries

\glsaddall % quitar luego, imprime todo el glosario diciendo que está en la primera página.

\begin{document}
	\frontmatter\pagestyle{empty} % marca el inicio de lo que ya no es portada, las primeras páginas en ser numeradas, aunque de forma distinta a las del mainmatter


	\pagestyle{plain}\setcounter{page}{1} %páginas comienzan a contabilizar en 1 (para evitar problemas con las portadas)
	\tableofcontents
	
	\clearpage
	\phantomsection
	\addcontentsline{toc}{chapter}{Resumen}
	
	\section*{Resumen}
	
	Estudiaremos los conceptos básicos de las bases de Schauder, circunscritas en el ámbito del Análisis Funcional. Veremos algunas de sus propiedades, tipos y utilidades en la caracterización de atributos de los espacios, centrándonos sobre todo en la reflexividad y en la estructura interna de los espacios, dando especial importancia a los espacios de sucesiones como $c_0$ y los $\ell_p$, destacando en especial a $\ell_1$ entre estos últimos. También daremos algunas notas relacionadas con la existencia y unicidad de bases de Schauder e intentaremos mostrar como estas bases son una extensión natural de las bases de Hamel en espacios finitos y una generalización de las bases de Hilbert a espacios de Banach más genéricos.
	
	\vspace*{2cm}
	
	{\let\clearpage\relax\section*{Abstract}}
	
	We will study the basic concepts of Schauder basis, mainly related to the Functional Analysis. We will see some of their properties, types and utilities in characterization of spaces attributes, focusing on reflexivity and the internal structure of spaces, giving special treat to the spaces of sequences like $c_0$ and $\ell_p$, laying emphasis on $\ell_1$ among latter. Furthermore, we will give some notes related to existence and uniqueness of Schauder basis and try to show how this basis are a natural extension of Hamel basis of finite spaces as well as a generalization  of Hilbert basis to more generic Banach spaces.
	
	\clearpage
	
	
	% para añdadir algo el índice: \addcontentline{toc}{tipo}{nombre}
	{\renewcommand{\chaptertitlesize}{25pt}
	\chapter*{Introducción}\addcontentsline{toc}{chapter}{Introducción}}
	
	
	Desde el inicio de la matemática moderna, con Descartes, hemos tenido el objetivo de asignar a los objetos de nuestro estudio una secuencia de cifras con las que poder simplificar dichos objetos y embeberlos en el mundo de los números, tratándolos así con todo el poder teórico de las matemáticas. Dio inicio esta tentativa con las primeras coordenadas \textit{cartesianas} ---así designadas en honor a su promotor--- que movilizaron el motor de la producción matemática haciéndolo más fructífero en los últimos $500$ años que en los previos $3\,300$ que hay desde las primeras ternas pitagóricas mesopotámicas conocidas.
	
	Es con esa filosofía en mente que surge la idea de realizar tal tentativa no ya con objetos físicos o con coordenadas finitas, sino con el poder del infinito. El objetivo es ahora el estudio de objetos mucho más abstractos: las funciones ---lo cual trajo consigo el problema de las inifitas dimensiones---.
	
	Al principio, el estudio se reducía a las funciones de $\mathcal{C}^\infty$ o infinitamente diferenciables ---las conocidas como series de Taylor, 1715, que fallaron en generar todo el espacio y se reducían a las funciones analíticas---. Con el tiempo, la idea fue madurando y se quiso ampliar el domino de validez a todas las funciones continuas, aún cuando no se tenía muy claro qué era una función continua todavía ---de aquí surgieron las series de Fourier, conocidas ya a principios del siglo XIX y que resultarían en una base no para las funciones continuas con la norma de la convergencia uniforme, sino para el espacio $L^2$ con la norma \emph{euclidiana} $\|\!\cdot\!\|_2$---. 
	
	Con el final del siglo XIX y la llegada del XX comenzaron a darse los mayores avances en este campo. Hilbert, Schmidt y Riesz---entre muchos otros--- comenzaron el estudio del espacio de las funciones de cuadrado integrable y consiguieron extraer ciertas nociones geométricas muy similares a las de los espacios euclidianos de dimensión finita, como por ejemplo la ortogonalidad. El concepto creció y se estableció la abstracción de los espacios de Hilbert ---término acuñado por von Neumann---, cuyo principal cometido y éxito es generalizar los espacios euclidianos a dimensión infinita respetando la existencia de sistemas ortonormales que generan la totalidad del espacio. Surgió así el estudio sistemático de la noción de base en dimensión infinita. 
	
	De forma paralela, Haar construyó un sistema ---el conocido como sistema de Haar--- de funciones que sirven de base no sólo para el espacio $L^2$, que ya tenía una base dada por las series de Fourier que además funcionaba por ser este un espacio de Hilbert, sino de todos los espacios $L^p$, con $1\leq p<\infty$ que  no son Hilbertianos. Poco después, Georg Faber demostró la existencia de una base para el espacio de las funciones continuas con la norma uniforme. 
	
	Todo este impulso histórico se cristalizó en la tesis de $1923$ de un joven matemático polaco, Juliusz Schauder. Él dio rigor y formalismo a la idea de una generalización de todo lo anterior introduciendo el concepto hoy en día conocido, en su honor, como \emph{bases de Schauder}: un avance más en el entendimiento de la estructura de los espacios de funciones inspirado en el objetivo de Descartes de simplificar el estudio de cualquier ente a un sistema de coordenadas coherente del que se puedan extraer conclusiones. 
	
	Es por querer dar valor a esta idea tan matemática que en el presente Trabajo Fin de Grado queremos dar un repaso a algunos de los resultados que las bases de Schauder han ayudado a probar y comprender.
	
	En el primer capítulo nos centraremos en la definición de lo que es una base de Schauder y caracterizaremos su existencia; también incluiremos algunos comentarios sobre el problema de la existencia de base para un espacio de Banach.
	
	En el Capítulo 2 estudiaremos la relación entre las bases de Schauder y el espacio dual, demostrando el primer gran resultado del trabajo: la caracterización de aquellos espacios de Banach que son reflexivos y admiten base de Schauder.
	
	En el Capítulo 3 estudiaremos las conocidas como sucesiones básicas, que no son más que bases de Schauder del espacio ---quizás menor que el total--- que generan. También hablaremos sobre la noción equivalencia para bases de Schauder.
	
	En el Capítulo 4 estudiaremos aquellas bases de Schauder que son incondicionales; es decir, es indiferente el orden que consideremos a la hora de estudiar la convergencia de la serie que define a cada vector del espacio. Caracterizaremos los espacios de Banach reflexivos con base incondicional.
	
	En el Capítulo 5 daremos por finalizada nuestra tentativa empleando la teoría de las bases de Schauder desarrollada a lo largo del trabajo para estudiar la estructura interna de los espacios de sucesiones $\ell_p$ con $1\leq p<\infty$, con el objetivo de dar a entender como de distintos son entre sí. 
	
	A mayores, queremos dejar patente que las bases de Schauder no son en absoluto una construcción artificiosa, sino que generalizan muy naturalmente a las bases de Hamel de los espacios de dimensión finita. Es por ello que iremos comparando sus propiedades en algunos capítulos para dejar claro el origen algebraico que tiene detrás esta idea y que tan bien resume la filosofía del Análisis Funcional de fusionar conceptos algebraicos con nociones topológicas para el estudio de las funciones.
	
	\begin{flushright}
		Santiago de Compostela, julio de 2025.
	\end{flushright}
	

	\mainmatter\pagestyle{fancy}\setlength{\parskip}{1.3mm}
	
	\chapter{Introducción a las Bases de Schauder}
	
	En este capítulo presentaremos algunos conceptos básicos: introduciremos la noción de base de Schauder para un espacio vectorial normado, caracterizaremos su construcción y existencia y más adelante, ya en la segunda sección del capítulo, veremos que la hipótesis de completitud es fundamental para obtener los resultados más importantes de este Trabajo Fin de Grado.
	
	También hablaremos, aunque de manera más superficial, del \textit{problema de existencia de base de Schauder} para espacios de Banach separables.
	
	\section{Bases de Schauder en espacios vectoriales normados}
	
	Es bien conocido que todo espacio vectorial tiene base de Hamel. Es decir, en todo espacio vectorial podemos encontrar un conjunto de vectores linealmente independientes cuyas combinaciones lineales (¡finitas!) generan el espacio. Además, todas las bases de Hamel de un mismo espacio vectorial tienen algo en común: su cardinal. Por tanto, es posible definir la \emph{dimensión} de cualquier espacio vectorial.
	
	No obstante, en el caso de dimensión infinita, la exigencia de finitud de las combinaciones lineales provoca que las bases de Hamel pierdan gran parte de su utilidad. Más exactamente (véase \cite[pág.48]{BST}):
	
	\begin{theorem}
		Ningún espacio vectorial de dimensión infinita admite base de Hamel numerable. 
	\end{theorem}
	
	Evidentemente, la no numerabilidad de una base es un gran inconveniente, pues reduce enormemente la capacidad de trabajar con éstas: con una base nosotros queremos reducir el número de elementos que necesitamos conocer de un espacio para poder trabajar con él, pero si dicha base es no numerable entonces es necesario conocer casi tantos elementos como tiene el espacio, lo cual hace redundante el sentido de una base.
	
	Para superar la problemática anteriormente mencionada podemos buscar ayuda en la topología y emplear el concepto de convergencia para generalizar las combinaciones lineales finitas a sumas infinitas, es decir, a series.
	
	Sea $(X,\|\!\cdot\!\|)$ un espacio vectorial normado.
	
	\begin{definition}[Base de Schauder]\label{def:bs} 
		Se dice que una sucesión $(e_i)_i\subset X$ es una \emph{base de Schauder de $X$} si para cada vector $x\in X$ existe una única sucesión de escalares $(a_i)_i$ de modo que 
		\[x=\sum_{i=1}^{\infty}{a_ie_i},\]
		es decir, cumpliendo que 
		\[\lim_{n\rightarrow\infty}\big\|x-\sum_{i=1}^{n}{a_ie_i}\big\|=0.\]
	\end{definition}
	
	La anterior definición ya nos deja entrever esa similitud de las bases de Schauder con las de Hamel: un conjunto de vectores---linealmente independientes--- que generan el espacio y nos permiten asignar de forma única unos coeficientes a todos los elementos de dicho espacio. Hasta aquí tenemos la filosofía algebraica del Análisis Funcional. 
	
	El lector podrá notar ahora los beneficios de introducir una perspectiva analítica, o más bien topológica, al poder reducir el tamaño de las bases. Donde antes teníamos bases no numerables para espacios de dimensión infinita, ahora tenemos una sucesión numerable que genera todo el espacio manteniendo la unicidad de los coeficientes (que tienen entonces derecho a ser denominados \emph{coordenadas}). Parece que vamos por buen camino en el cumplimiento de nuestro objetivo.
	
	Veamos ahora una caracterización de estas bases que nos será muy útil tanto para demostrar que una sucesión es, en efecto, base de Schauder como para construirlas.
	
	\begin{definition}[Proyecciones canónicas de una base de Schauder]
		Si $(e_i)_i$ es una base de Schauder de $X$, la sucesión de \emph{proyecciones canónicas} asociada a $(e_i)_i$ es la sucesión de operadores $(P_n)_n\subset\mathcal{B}(X)$ dada, para cada $n\in\mathbb{N}$, por
		\begin{equation}\label{eq:1}
			P_n\colon x=\sum_{i=1}^\infty a_ie_i\in X\longmapsto P_nx=\sum_{i=1}^n a_ie_i\in X.
		\end{equation}
	\end{definition}
	
	%{\color{red} Texto}
	
	Sea $(X,\|\!\cdot\!\|)$ un espacio vectorial normado con base de Schauder $(e_i)_i$.
	
		
	\begin{proposition}[Caracterización de las proyecciones canónicas]\label{proyections}
		La sucesión de proyecciones canónicas asociada a $(e_i)_i$ satisface las siguientes propiedades:
		\begin{itemize}
			\item[$(a)$]{para cada $n\in\mathbb{N}$, $P_n(X)$ es un espacio vectorial de dimensión $n$;}
			\item[$(b)$]{para cualesquiera $n,\,m\in\mathbb{N}$, $P_nP_m=P_mP_n=P_{\min\{n,m\}}$;}
			\item[$(c)$]{para cada $x\in X$, $\|P_nx-x\|\to 0$ cuando $n\to\infty$.}
		\end{itemize}
		Además, dada una sucesión de operadores $(P_n)_n\subset\mathcal{B}(X)$ satisfaciendo $(a)$, $(b)$ y $(c)$, existe una base de Schauder de $X$ que tiene a $(P_n)_n$ como sucesión de proyecciones canónicas asociada.
	\end{proposition}
	
%		\begin{remark*}
		\begin{nota}
			Antes de ver la demostración de la Proposición \ref{proyections} conviene recordar que se dice que un operador $P\in\mathcal{B}(X)$ es una \emph{proyección} (lineal y continua) si $P^2\equiv PP=P$. En tal caso, $I-P$ también es una proyección. Además, $P(X)={\rm ker}\,(I-P)$, luego $P$ restringido a $P(X)$ es la identidad. Análogamente, ${\rm ker}\,P=(I-P)(X)$. Es evidente que si $P$ es una proyección no trivial ($P\neq 0$ y $P\neq I$), entonces $\|P\|\geq 1$ y $\|I-P\|\geq 1$.
%		\end{remark*}
	\end{nota}		
		
		\begin{proof}
			\noindent``$\Longrightarrow$''
			\begin{itemize}
				\item[$(a)$]{A la vista de (\ref{eq:1}) es evidente que ${\rm dim}\,P_n(X)\leq n$. Además, $\{e_i\colon i\in\mathbb{N}\}$ es un conjunto linealmente independiente, pues si $\lambda_1e_1+\lambda_2 e_2+\cdots+\lambda_ne_n=0$, dado que 
					\[0=\sum_{i=1}^{\infty}0\,e_i=\sum_{i=1}^n{\lambda_ie_i} + \sum_{i=n+1}^{\infty}0\,e_i,\]
					la unicidad de los escalares que se exige en la definición de base de Schauder implica que $\lambda_i=0$ para todo $1\leq i\leq n$.  Por tanto, $\dim\,P_n(X)=n$.}
				\item[$(b)$]{A la vista de (\ref{eq:1}) es evidente que $P_nP_m=P_mP_n=P_{\min\{n,m\}}$.}
				\item[$(c)$]{Dado que $(P_n(x))_n$ es la sucesión de sumas parciales de la serie (convergente) que coincide con $x$, por definición, $\|P_n(x)-x\|\to 0$ cuando $n\to\infty$.}
			\end{itemize}\medskip
			
			\noindent``$\Longleftarrow$''
			
			Indicamos a continuación las ideas claves de la prueba. Se detallan todos los pasos de manera más pormenorizada al final de este epígrafe (en la página \pageref{D:proyecciones}).
			
			Sea $P_0\equiv 0$ la aplicación nula y consideremos, para cada $i\in I$, un vector $e_i\in P_i(X)\cap 
			{\rm ker}\,P_{i-1}$. En tal caso, para cada $x\in X$ se tiene que
			\[x=\lim_{n\to\infty}{P_n(x)}=\lim_{n\to\infty}\sum_{i=1}^nP_i(x) - P_{i-1}(x)=  \sum_{i=1}^\infty a_ie_i,\]
			ya que $\dim\,P_i(X)/P_{i-1}(X)=1$. La unicidad de la sucesión de escalares $(a_i)_i$ exigida en la definición de base de Schauder es consecuencia de que $a_ie_i=P_i(x)-P_{i-1}(x)$ para cada $i\in\mathbb{N}$.
		\end{proof}
	
	
	El lector podrá notar que la idea detrás de esta caracterización es muy interesante ---y bella---: para construir las bases del espacio total se parte de un espacio de dimensión finita y se va haciendo más grande hasta que, en una especie de convergencia espacial, se construye una sucesión de espacios que se aproxima cada vez más al espacio total. Nos gustaría incidir en esta idea porque guarda una estrecha relación con la filosofía del Análisis Matemático: estudiar aquello que no comprendemos fraccionándolo primero en aquello que sí conocemos y extrapolando luego mediante el uso de límites sus propiedades.
	
	Además, remarca de nuevo la idea principal de este Trabajo Fin de Grado: las bases de Schauder pretenden ser bases de Hamel que admiten sumas infinitas, es decir, \emph{límites} de bases de Hamel.
	
	Veamos ahora un resultado que abre las puertas a lo que será nuestro ámbito de estudio: los espacios completos.
	
	\begin{proposition}[Condición suficiente de base de Schauder de la compleción $\hat{X}$ de $X$]\label{P:1.4}	
		Si
		\begin{equation}\label{eq:acotacion}
			\sup_{n\in\mathbb{N}}\|P_n\|<\infty
		\end{equation}
		la sucesión $(e_i)_i$ también es una base de Schauder de la compleción $\hat{X}$ de $X$.
	\end{proposition}
	
		\begin{proof}			
			Consideremos, para cada $n\in\mathbb{N}$, la extensión continua de $P_n$ a $\hat{X}$; es decir, para cada $\hat{x}\in\hat{X}$, definimos $\hat{P}_n(\hat{x})= \lim_{k\to\infty}P_n(x_k)$, siendo $(x_k)_k\subset X$ cualquier\footnote{Las proyecciones canónicas $P_n$ son continuas. Por tanto, si $(x_k)_k$ e $(y_k)_k$ son dos sucesiones convergentes a $\hat{x}$, entonces $\hat{P}_n(x_k-y_k)\to 0$ cuando $k\to\infty$ y equivalentemente $\lim_{k\to\infty}P_n(x_k)=\lim_{k\to\infty}P_n(y_k)$.} sucesión convergente a $\hat{x}$.
			
			Veamos que se cumplen las propiedades $(a)$, $(b)$ y $(c)$ del enunciado del Teorema \ref{proyections}.
			
			\begin{itemize}
				\item[$(a)$]{Atendiendo a la definición de $\hat{P}_n$, resulta evidente que 
					$\hat{P}_n(\hat{X})\subset\overline{P_n(X)}$. Además, como ---por hipótesis--- $P_n(X)$ es un subespacio vectorial de dimensión finita, en particular es cerrado, luego $P_n(X)=\overline{P_n(X)}$. Por tanto $\hat{P}_n(\hat{X})\subset P_n(X)$.
					
					Por otra parte, dado que $P_n(X)=\hat{P}_n(X)\subset\hat{P}_n(\hat{X})$, concluimos finalmente que $P_n(\hat{X})$ es un espacio vectorial de dimensión $n$, ya que ---por hipótesis--- ${\rm dim}\,P_n(X)=n$.}
				
				\item[$(b)$]{Dado $\hat{x}\in \hat{X}$ arbitrario pero fijado, existe $(x_k)_k\subset X$ tal que $x_k\to\hat{x}$ cuando $k\to\infty$. Así pues, para $m$, $n\in\mathbb{N}$ cualesquiera,
					\begin{align*}
						\hat{P}_n\hat{P}_mx & = \lim_{k \to\infty}{\hat{P}_n\hat{P}_mx_k} = \lim_{k \rightarrow\infty}{\hat{P}_nP_mx_k} = \lim_{k \rightarrow\infty}{P_nP_mx_k}\\ & =  \lim_{k \rightarrow\infty}{P_mP_nx_k} = \lim_{k \rightarrow\infty}{\hat{P}_m\hat{P}_nx_k} = \hat{P}_m\hat{P}_nx,
					\end{align*}
					por lo que $\hat{P}_n\hat{P}_m=\hat{P}_m\hat{P}_n$. La igualdad $\hat{P}_n\hat{P}_m=\hat{P}_{\min\{n,m\}}$} es consecuencia evidente de lo anterior y de la hipótesis $P_nP_m=P_{\min\{n,m\}}$.
				\item[$(c)$]{Conviene observar inicialmente que, para cada $n\in\mathbb{N}$, se tiene que $\|\hat{P}_n\|\leq\|P_n\|$. En efecto, pues para cualquier $\hat{x}\in\hat{X}$ se tiene que
					\[\|\hat{P}_n(\hat{x})\| = \|\lim_{k \to \infty}P_n(x_k)\| = \lim_{k\to \infty}\|P_n(x_k)\| \leq \lim_{k \to \infty}\|P_n\|\|x_k\| = \|P_n\|\|\hat{x}\|,\]
					siendo $(x_k)_k\subset X$ una sucesión tal que $x_k\to \hat{x}$ cuando $k\to\infty$.
					
					Así pues, dado $\hat{x}\in\hat{X}$ arbitrario pero fijado, considerando $(x_k)_k\subset X$ tal que $x_k\to \hat{x}$ cuando $k\to\infty$, tenemos que
					\begin{align*}   
						\|\hat{x}-\hat{P}_n\hat{x}\|&=\|\hat{x}-x_{k}+x_{k}-\hat{P}_nx_{k}+\hat{P}_nx_{k}-\hat{P}_n\hat{x}\|\\ 
						&\leq\|\hat{x}-x_{k}\|+\|x_{k}-P_nx_{k}\|+\|\hat{P}_nx_{k}-\hat{P}_n\hat{x}\|\\ 
						&\leq\|\hat{x}-x_{k}\|+\|x_{k}-P_nx_{k}\|+\|\hat{P}_n\|\,\|x_{k}-\hat{x}\|\\
						&\leq\|\hat{x}-x_{k}\|+\|x_{k}-P_nx_{k}\|+\sup_{n\in\mathbb{N}}\|\hat{P}_n\|\,\|x_{k}-\hat{x}\|\\
						&\leq \|\hat{x}-x_{k}\|+\|x_{k}-P_nx_{k}\|+\sup_{n\in\mathbb{N}}\|P_n\|\,\|x_{k}-\hat{x}\|.
					\end{align*}
					La condición $x_k\to \hat{x}$ cuando $k\to\infty$ asegura la existencia de $k_\varepsilon\in\mathbb{N}$ de modo que si $k\geq k_\varepsilon$, entonces
					\[\|\hat{x}-x_{k}\|<\dfrac{\varepsilon}{3}\quad\text{y}\quad \sup_{n\in\mathbb{N}}\|P_n\|\,\|x_{k}-\hat{x}\|<\dfrac{\varepsilon}{3}.\]
					Por otra parte, dado que $(P_n)_n$ es una sucesión de proyecciones canónicas, las sucesión $(P_n)_n$ satisface la propiedad $(b)$ del enunciado de la Proposición \ref{proyections}, luego existe $N_{k_\varepsilon}\in\mathbb{N}$ de modo que si $n\geq N_{k_\varepsilon}$ entonces
					\[\|x_{k}-P_nx_{k}\|<\dfrac{\varepsilon}{3}.\]
					
					Por tanto, $\|\hat{x}-\hat{P}_n\hat{x}\|<\varepsilon$ si $n\geq N_{k_\varepsilon}$; es decir, $\|\hat{x}-\hat{P}_n\hat{x}\|\to 0$ cuando $n\to\infty$.\qedhere}
			\end{itemize}
		\end{proof}
	
%	\begin{remarkimp*}
	\begin{nota}
		La condición (\ref{eq:acotacion}) sugiere definir
		\[\vertiii{x}=\sup_n\|P_nx\|\quad\text{para cada $x\in X$.}\]
		Comprobar que $\vertiii{\cdot}$ define una norma en $X$ es muy sencillo:
		\begin{itemize}
			\item[(N1)]{$\vertiii{x}\geq 0$ trivialmente y $0=\vertiii{x}=\sup_{n}{\|P_nx\|}$ implica que  $\|P_nx\|=0$ para todo $n\in\N$, pero como $\|\!\cdot\!\|$ es una norma, podemos deducir que $P_nx=0$ para todo $n$, de donde se sigue finalmente que $x=0$ por ser $(e_i)_i$ una base de Schauder.}
			\item[(N2)]{$\vertiii{\lambda x}=\sup_n\|P_n(\lambda x)\|=\sup_n|\lambda|\,\|P_nx\|=|\lambda|\, \sup_n\|P_nx\|=\lambda\,\vertiii{x}$.}
			\item[(N3)]{dado que $\|P_nx+P_ny\|\leq \|P_n\|+\|P_ny\|$,
				\begin{align*}
					\vertiii{x+y}&=\sup_n\|P_n(x+y)\|=\sup_n\|P_nx + P_ny\|\leq\sup_n(\|P_nx\|+\|P_ny\|)\\ 
					&\leq\sup_n\|P_nx\|+\sup_n\|P_ny\|= \vertiii{x} + \vertiii{y}.
			\end{align*}}
		\end{itemize}
%	\end{remarkimp*}
	\end{nota}
	
	\section{Bases de Schauder en espacios de Banach}
	
	Añadamos pues la hipótesis de completitud de la que jamás nos desharemos, pues resulta imprescindible para la acotación uniforme de las proyecciones, hecho que ---tal y como el lector verá con el paso de las páginas---, asegura la convergencia de series, continuidad de aplicaciones  y existencia de espacios que describiremos en los sucesivos capítulos.
	
	Sea $(X,\|\!\cdot\!\|)$ un espacio de Banach con base de Schauder $(e_i)_i$.
	
	\begin{proposition}[Auxiliar para el Teorema \ref{th:banach} siguiente]\label{P:1.5}
		La sucesión $(e_i)_i$ es base de Schauder de $(X,\vertiii{\cdot})$.
	\end{proposition}
	
		\begin{proof}
			Probaremos que la sucesión de proyecciones canónicas $(P_n)_n$ satisface las condiciones enunciadas en la Proposición \ref{proyections}.
			
			Las propiedades $(a)$ y $(b)$ no dependen de la norma, por lo que se satisfacen automáticamente. Por otra parte, la propiedad $(c)$ es sencilla de probar, pues para cada $x\in X$ se tiene que
			\begin{align*}\label{eq:unicidad}  
				\vertiii{x-P_nx} &=\sup_{m}{\|P_m(x-P_nx)\|}=\sup_{m}{\|P_mx-P_mP_nx\|}=\sup_{m}{\|P_mx - P_nP_mx\|}\\&=\sup_{n\leq m}{\|P_mx - P_nP_mx\|}=\sup_{n\leq m}{\|P_mx - P_nx\|}\to 0\quad\text{cuando $n\to\infty$}  
			\end{align*}
			porque, para cualquier $y\in X$, por ser $(e_i)_i$ base de Schauder de $(X,\|\!\cdot\!\|)$, tenemos que $P_ny \to y$ cuando $n\to\infty$ en $(X,\|\!\cdot\!\|)$.
		\end{proof}
	
	
	\begin{theorem}[Stefan Banach: «Completitud $\Longrightarrow$ Acotación de las proyecciones canónicas»]\label{th:banach}
		La sucesión de proyecciones canónicas $(P_n)_n$ es (uniformemente) acotada en $(\mathcal{B}(X),\|\!\cdot\!\|)$.
		
	\end{theorem}
	

	\begin{proof}
		Conviene observar que
		\[\vertiii{P_n}=\sup_{\vertiii{x} \leq 1}{\vertiii{P_nx}}=\sup\,\Big\{\sup_{m \leq n}\|P_mx\| \colon \vertiii{x} = \sup_{k}\| P_kx\| \leq 1\Big\} \leq 1,\]
		por lo que la sucesión $(P_n)_n$ es (uniformemente) acotada en $\mathcal{B}(X)$ si consideramos a $X$ dotado de la norma $\vertiii{\cdot}$ y basta entonces demostrar que $\vertiii{\cdot}$ es equivalente a $\|\!\cdot\!\|$, es decir, que existen constantes $M_1$, $M_2>0$ de modo que $\vertiii{x}\leq M_1\|x\|\leq M_2\vertiii{x}$ para todo $x\in X$.
		
		En efecto, pues tal caso, la acotación uniforme en $(X,\|\!\cdot\!\|)$ será inmediata:
		\[\|P_n\|=\sup_{\|x\|\leq 1}\|P_n x\|\leq\sup_{\vertiii{x}\leq M_1}M_2\vertiii{P_nx}=\dfrac{M_2}{M_1}\sup_{\vertiii{x}\leq 1}\vertiii{P_nx}\leq\dfrac{M_2}{M_1}\quad\text{para todo $n\in\mathbb{N}$}.\]
		
		
		Si $\left(X,\vertiii{\cdot}\right)$ fuese completo, dado que la aplicación $I \colon \left(X, \vertiii{\cdot}\right)\longrightarrow (X, \|\!\cdot\!\|)$ es trivialmente continua (pues $\|x\|\leq\vertiii{x}$ para todo $x\in X$) y biyectiva, el { Teorema \ref{th:ap-abierta} de la aplicación abierta} nos permitiría afirmaría que $I^{-1}$ también es continua, por lo que $\vertiii{\cdot}$ y $\|\!\cdot\!\|$ serían necesariamente normas equivalentes.
		
		Probaremos que es posible proceder de tal modo. Para ello, demostraremos que $\left(X,\vertiii{\cdot}\right)$ es un espacio de Banach mostrando que $\hat{X}\subset X$, siendo ahora $\hat{X}$ la compleción de $X$ respecto de la norma $\vertiii{\cdot}$.
		
		Dado $x\in\hat{X}$, como $(e_i)_i$ es base de Schauder de $\hat{X}$ (esto es una consecuencia de la Proposición \ref{P:1.4} y la Proposición \ref{P:1.5}), existe una única sucesión de escalares $(\alpha_i)_i$ de modo que
		\[x=\lim_{n\to\infty}\sum_{i=1}^n \alpha_i e_i\quad\text{en $(\hat{X},\vertiii{\cdot})$;}\]
		luego $(\alpha_1e_1+\cdots+\alpha_n e_n)_n$ es una sucesión de Cauchy en $(X,\vertiii{\cdot})$, y dado que $\|\!\cdot\!\|\leq\vertiii{\cdot}$, también es una sucesión de Cauchy en $(X,\|\!\cdot\!\|)$. Por tanto, por ser $(X,\|\!\cdot\!\|)$ completo, existe $x'\in X$ de modo que
		\[x'=\lim_{n\to\infty}\sum_{i=1}^n \alpha_i e_i\quad\text{en $(X,\|\!\cdot\!\|)$.}\]
		Es decir, $x'=x\in X$.
	\end{proof}
	
	En vista de este teorema podemos introducir la siguiente definición:
	
	\begin{definition}[Constante básica]
		El número real $\text{bc}\,(e_i)_i:=\sup_{n\in\N}\|P_n\|$ recibe el nombre de \emph{constante básica} de $(e_i)_i$.
	\end{definition}
	
	Como pequeño complemento a la definición anterior, tenemos el concepto de base \textit{monótona}, que es aquella para la que $\text{bc}\,(e_i)_i = 1$. La denominación es explicativa, pues refleja que en dicha base, para cualquier selección de escalares $(a_i)_i$, la sucesión de números reales
	\[\Big(\big\|\sum_{i=1}^na_ie_i\big\|\Big)_n\]
	es monótona creciente, lo cual es evidente pues, tomando $n>m$,
	\[\big\|\sum_{i=1}^m a_ie_i\big\|=\|P_mx\|=\|P_nP_mx\|\leq \|P_nx\|\|P_m\|\leq \|P_nx\|= \big\|\sum_{i=1}^n a_ie_i\big\|.\]
	
	
	\begin{center}
		\begin{minipage}{0.35\textwidth}\centering\small
%			\includegraphics[height=4cm]{figuras/schauder.png}\\[2mm]
			\href{https://mathshistory.st-andrews.ac.uk/Biographies/Schauder/}{Juliusz Schauder}\\
			Leópolis, 1899 – Leópolis, 1943
		\end{minipage}
		\begin{minipage}{0.60\textwidth}
			
			Schauder fue un matemático polaco--judío con una corta pero relevante trayectoria investigadora en los campos del análisis funcional, EDP's y física matemática.\smallskip
			
			Participó en la Primera Guerra Mundial. Fue alumno de Banach y Steinhaus en Leópolis, donde se doctoró en $1923$.\smallskip
			
			Colaboró con eminentes matemáticos, como Bernstein, Sobolev, Lewy o Leray, con quien llegó incluso a compartir un premio.
		\end{minipage}
	\end{center}\medskip
	
	Para finalizar esta sección incluimos una breve mención al conocido problema de la existencia de base.
	
	Con el nacimiento de la idea de las bases de Schauder surgió también una cuestión natural:
%	\centerline{¿todo espacio de Banach separable admite base de Schauder? }
	
	Es una pregunta que surge del teorema con el que dimos inicio a este capítulo. Tendríamos la esperanza de que la respuesta fuese afirmativa y entonces podríamos añadir otra analogía entre las bases de Hamel y las de Schauder. Pero el interrogante se aclaró en $1973$, cuando el matemático sueco Per Enflo \cite{Enflo1973} proporcionó un ejemplo de espacio de Banach separable que no admite base de Schauder. En cualquier caso, todos los espacios clásicos sí admiten bases de Schauder ---como el lector podrá ver en el Capítulo \ref{C:ejemplos} de este Trabajo Fin de Grado---.
	
	Por otra parte, aún cabe la posibilidad de que todo espacio de Banach separable tenga algún subespacio con base de Schauder (ignorando evidentemente por triviales a los subespacios de dimensión finita). El lector puede esperar al Capítulo \ref{C:sucesiones_basicas} donde se ahondará más en esta cuestión.
	
	
%	\renewcommand*{\proofname}
	\begin{proof}[Demostración detallada de la Proposición \ref{proyections}]\label{D:proyecciones}
		
		\noindent\emph{Existencia de los vectores $e_i$ satisfaciendo las condiciones requeridas}.
		
		\begin{addmargin}[5mm]{0mm}
			En virtud de $(b)$ sabemos que
			\[P_{\min\{n,m\}}(X)=P_nP_m(X) \subset P_n(X)\quad\text{y}\quad P_{\min\{n,m\}}(X)=P_mP_n(X)\subset P_m(X).\]
			Por tanto, $P_{\min\{n,m\}}(X) \subset P_{\max\{n,m\}}(X)$. Es decir,
			\[P_0(X) \subset P_1(X) \subset \cdots \subset P_n(X) \subset P_{n+1}(X) \subset P_{n+2}(X) \subset \cdots\]
			Por otra parte, en virtud de $(a)$, $\dim\,P_n(X)-\dim\,P_{n-1}(X)=1$. Por tanto, para cada $n\in\mathbb{N}$, existe un vector no nulo $\hat{e}_n\in P_n(X)\setminus P_{n-1}(X)$. Trivialmente, $\{\hat{e}_i\colon i\in\mathbb{N}\}$ es un conjunto linealmente independiente.
			
			Además, como $X=P_{n-1}(X)\oplus{\ker P_{n-1}}$, tenemos en realidad que $\hat{e}_n\in P_n(X)\cap {\rm ker}P_{n-1}$.
			
			Sea ahora, para cada $i\in\mathbb{N}$, $e_i=P_i\hat{e_i}$. En tal caso, por la propiedad $(b)$, dado un índice $j>i$, se tiene que $P_je_i=P_jP_i\hat{e}_i=P_i\hat{e}_i=e_i$. Es decir,
			\begin{equation}\label{eq:3}
				P_je_i=e_i\quad\text{para todo $j>i$.}    
			\end{equation}
			Por otra parte, si $j<i$, nuevamente por la propiedad $(b)$, $P_je_i=P_jP_1\hat{e}_i=P_j0=0$. Es decir,
			\begin{equation}\label{eq:4}
				P_je_i=0\quad\text{para todo $j<i$.}    
			\end{equation}
			Además, ahora es sencillo comprobar que $\{e_i\colon i\in\mathbb{N}\}$ sigue siendo un conjunto linealmente independiente.
		\end{addmargin}
		
		\noindent\emph{Para cada $x\in X$, los coeficientes $(a_i)_i$ son únicos.}
		\begin{addmargin}[5mm]{0mm}
			En principio, para cada $k\in\mathbb{N}$, puesto que $P_kx\in{\rm span}\,\{e_i\, \colon 1\leq i\leq k\}$, sabemos que existen escalares $a^k_i$ con $1\leq i\leq k$ tales que $P_kx=a_1^ke_1+\cdots +a_k^k e_n$. 
			
			Ahora bien, en realidad, $a_i^m=a_i^n\equiv a_i$ para cualesquiera $m$, $n$ e $i\in\mathbb{N}$. En efecto, pues en virtud de (\ref{eq:3}), la linealidad de las proyecciones, (\ref{eq:4}) y la propiedad $(b)$, tenemos que (si $n<m$):
			\begin{align*}
				\sum_{i=1}^na_i^m e_i&=\sum_{i=1}^na_i^mP_ne_i=P_n\Big(\sum_{i=1}^n a_i^me_i\Big)+P_n\Big(\sum_{i=n+1}^m a_i^me_i\Big)=P_n\Big(\sum_{i=1}^m a_i^me_i\Big)\\
				&=P_nP_mx=P_mP_nx=P_m\Big(\sum_{i=1}^n a_i^ne_i\Big)=\sum_{i = 1}^n{a_i^nP_me_i}.
			\end{align*}
			Por tanto, dado que ---de nuevo en virtud de (\ref{eq:3})--- $P_me_i=e_i$ para todo $1\leq i\leq n(<m)$, la independencia lineal del conjunto $\{e_i\colon i\in\mathbb{N}\}$ asegura la coincidencia de las coordenadas. Así pues,
			\[P_ix-P_{i-1}x=\sum_{j=1}^i a^i_je_j-\sum_{j=1}^{i-1}a^{i-1}_je_j=a_ie_i.\qedhere\]
		\end{addmargin}
	\end{proof}
	
	\renewcommand*{\proofname}{Demostración}
	
	
	 %marca el inicio de la parte principal
	
	\chapter{Motivación e idea general}
	
	Tradicionalmente, consideramos en el contexto de la regresión que solo
	existe una tendencia que influencia a la variable de estudio, $Y$, en
	relación a las covariables $X$\footnote{Por ejemplo, en la regresión
		lineal en media.}. Puedes citar un enlace como
	\href{https://www.anthropic.com}{este} y se mostrará en negrita y en
	el color primario.
	
	Por tanto, surge un problema con estas técnicas cuando se observan diferentes tendencias causadas, por ejemplo, por variables explicativas no observadas. Imaginemos una situación en la que le damos un mismo examen de acceso a la universidad a todos los alumnos de un instituto. Nuestra variable respuesta es «Nota en un examen», y queremos estudiar la relación que existe con la variable «cantidad de horas estudiadas». Podríamos suponer a modo teórico que a un mayor número de horas estudiadas mayor será la nota en el examen, y que dicha relación será más o menos lineal. Es importante notar que nuestro estudio no está teniendo en cuenta ninguna variable que clasifique a cada individuo en una u otra clase.
	
	Ahora bien, en el instituto encontramos clases de diferentes años, y por ello van a contestar de forma distinta al mismo examen: los cursos más altos tendrán una ventaja con respecto al resto de alumnos, obteniendo mejores resultados que cursos más bajos. Si ahora realizamos un estudio de regresión en media obtendríamos una única predicción para cada unidad de tiempo de estudio, obviando esa estructura subyacente que desconocemos, pero que podríamos intentar detectar. 
	
	Simplifiquemos este ejemplo aún más, y consideremos la misma prueba pero añadiendo a los alumnos del colegio de educación primaria de al lado. Ahora, nuestros datos se van a acumular en dos grandes zonas bastante separadas: por un lado los alumnos de instituto que por lo menos son capaces de responder algunas de las preguntas, y los alumnos del colegio que no cuentan con los conocimientos para responder casi nada. Si realizásemos de nuevo una regresión en media, sucederá que la predicción caerá en una región de valores de $Y$ en la que no vamos a tener observaciones 
	
	\section{Ecuaciones}
	
	Esta ecuación se referencia más adelante, ver \eqref{eq:einstein}.
	\begin{equation}
		E = mc^2 \label{eq:einstein}
	\end{equation}
	
	Esta otra tiene etiqueta pero nunca se referencia, así que no debe
	llevar número:
	\[
	a^2+b^2=c^2 
	\]
	
	
	Como dijimos en \eqref{eq:einstein}, la masa y la energía están
	relacionadas.
	
	\section{Definiciones y teoremas}
	
	\begin{definition}{Espacio de Banach}
		Un espacio normado $(X,\|\cdot\|)$ es un \textbf{espacio de Banach} si
		es completo respecto a la métrica inducida por la norma.
	\end{definition}
	
	\begin{theorem}{Teorema de Hahn--Banach}
		Sea $X$ un espacio vectorial normado y $Y\subset X$ un subespacio.
		Todo funcional lineal acotado en $Y$ se puede extender a todo $X$ sin
		aumentar su norma.
	\end{theorem}
	
	\begin{proof}
		Aquí iría el desarrollo de la demostración, con la misma sangría y
		regla vertical izquierda que en la plantilla original de Typst.
		\begin{itemize}
			\item 1
		\end{itemize}
		
	\end{proof}
	
	
	\section{Figuras}
	
	\begin{figure}[h]
		\centering
		\fbox{\parbox{5cm}{\centering\vspace{1.5cm}Imagen de ejemplo\vspace{1.5cm}}}
		\caption{Ejemplo de figura con numeración coloreada.}
		\label{fig:ejemplo}
	\end{figure}
	
	Como se ve en la Figura~\ref{fig:ejemplo}, el pie de figura también
	sale coloreado.
	
	\chapter{Segundo capítulo}
	
	\section{Otra sección}
	
	Texto de relleno para comprobar que la cabecera reaparece en las
	páginas siguientes y desaparece en la primera página de cada capítulo.
	
	\begin{nota}
		A ver, alguien puede decirme cuanto es la suma entre la hipotenusa de dos catetos y los catetos de las dos serpientes que atacan el puerto de san fernando de osiris?
		
		Por tanto, surge un problema con estas técnicas cuando se observan diferentes tendencias causadas, por ejemplo, por variables explicativas no observadas. Imaginemos una situación en la que le damos un mismo examen de acceso a la universidad a todos los alumnos de un instituto. Nuestra variable respuesta es «Nota en un examen», y queremos estudiar la relación que existe con la variable «cantidad de horas estudiadas». Podríamos suponer a modo teórico que a un mayor número de horas estudiadas mayor será la nota en el examen, y que dicha relación será más o menos lineal. Es importante notar que nuestro estudio no está teniendo en cuenta ninguna variable que clasifique a cada individuo en una u otra clase.
		
		\[
		1+1 = 2
		\]
		
		Ahora bien, en el instituto encontramos clases de diferentes años, y por ello van a contestar de forma distinta al mismo examen: los cursos más altos tendrán una ventaja con respecto al resto de alumnos, obteniendo mejores resultados que cursos más bajos. Si ahora realizamos un estudio de regresión en media obtendríamos una única predicción para cada unidad de tiempo de estudio, obviando esa estructura subyacente que desconocemos, pero que podríamos intentar detectar.
	\end{nota}
	
%	\verb|hola que tal estás| .
	
\begin{codigo}
cluster = 1

print(cluster)

# a ver cómo se ve una función
function abrir(){
    foo = seq(1,2,1:100)
    cat("que mierda pasa")
    pintar()

}
\end{codigo}
	
%	\renewcommand{\appendixtitle}{Anexos}
%	
%	% ----------------------------------------------------------------------
%	% Apéndices: \appendix hace que \thechapter pase de usar números (1,2,3)
%	% a usar letras (A,B,C) automáticamente — no hay que tocar nada más,
%	% el mismo \@makechapterhead ya lo pinta bien porque usa \thechapter.
%	% ----------------------------------------------------------------------
	\appendix
	
	\chapter{Resultados auxiliares}
	
	% Improtante recordar citar el libro de donde saco estes ejercicios.
	\begin{lemma} \label{lm:incondicionalidad-reordenacion}
		Sea $\sum{x_i}$ una serie en un espacio de Banach $X$ y $x \in X$. Son equivalentes:
		\begin{itemize}
			\item[$(i)$] Para todo $\varepsilon > 0$ existe un conjunto finito $F \subset \mathbb{N}$ tal que 
			$$\big\| x - \displaystyle \sum_{i \in F'}{x_i} \big\| < \varepsilon$$
			para cualquier $F'$ conjunto finito de $\mathbb{N}$ tal que $F \subset F'$
			\item[$(ii)$] Si $\pi$ es una permutación de $\mathbb{N}$ entonces 
			$$\sum{x_{\pi(i)}}=x$$
		\end{itemize}
	\end{lemma}
	
		\begin{proof}
			\noindent $(i) \Longrightarrow (ii)$
			
			Se verá que $F$ está contenido en el subconjunto de las permutaciones de los primeros $m$ números naturales y, aplicando la hipótesis tomando por $F'$ ese subconjunto, se tendrá el resultado.
			
			Sea $\varepsilon >0$ fijado y considérese el conjunto $F \subset \mathbb{N}$ como el que verifica la hipótesis. Como $\pi$ es una permutación de $\mathbb{N}$, y por tanto una biyección, y $F$ es un subconjunto finito de $\mathbb{N}$ se tiene que para cada $n \in F$ existe un $m$ tal que $\pi(m) = n$, y como $F$ es finito existe un cierto $n_0 \in \mathbb{N}$, que puede ser el máximo de los $m$, tal que $F \subset \{ \pi(1), \pi(2),\ldots, \pi(n_0) \}$. Por la hipótesis, se tiene que 
			$$\big\| x -  \sum_{i=1}^{n}{x_{\pi(i)}} \big\| < \varepsilon$$
			
			\noindent $(ii) \Longrightarrow (i)$
			
			Suponiendo que no sucede (ii) pero si (i), se construirá una sucesión de conjuntos finitos $F_k$ para los cuales la suma de sus elementos sea siempre mayor a un cierto valor, después se escogerá una permutación que deja dichos conjuntos fijos fijos y se verá que la serie asociada a esa permutación no converge.
			
			Por reducción al absurdo, supóngase  que existe un $\varepsilon > 0$ para el cual todo conjunto finito $F$ verifica que existe un $F' \supset F$ tal que  
			$$\big\| x -  \sum_{i \in F'}{x_i} \big\| > \varepsilon$$ 
			
			Por hipótesis, y como la identidad es una permutación de $\mathbb{N}$, $\sum{x_i} = x$ y por tanto para todo $k \in \mathbb{N}$ existe un $n'_k$ tal que 
			$$\big\| \displaystyle \sum_{i=1}^{n'_k}{x_i} - x \big\| < \frac{1}{k}$$.
			
			Sea ahora, para $k = 1$ un subconjunto $F' \subset \mathbb{N}$   que contiene a $\{1,\ldots, n'_1\}$ y que verifica que 
			$$\big\| x - \displaystyle \sum_{i \in F'}{x_i} \big\| > \varepsilon$$
			(y que existe por la hipóteis de reducción al absurdo). Sea ahora $F_1 = F' \cap \mathbb{N} \backslash \{1, \ldots, n'_0\}$, $n_1 = n'_1$ y $n_2 = \max\{n'_2, \max(F_1) +1\}$, así se tiene que  $n_{2} >  \max{F_1} \geq \min{F_1} > n_1$.
			
			Aplicando inducción, se construyen todos los conjuntos $F_k$ y los naturales $n_k$ verificando que $n_{k+1} > \max F_k \geq \min F_k > n_k$ y 
			$$\big\| \sum_{i \in F'}{x_i} - x \big\| = \big\|  \sum_{i=1}^{n_k}{x_i} +  \sum_{i \in F_k}{x_k} - x\big\| > \varepsilon$$.
			Por tanto 
			$$\varepsilon < \big\|\sum_{i=1}^{n_k}{x_i} + \sum_{i \in F_k}{x_i} - x \big\| \leq \big\|  \sum_{i=1}^{n_k}{x_i}  - x\big\| + \big\| \sum_{i \in F_k}{x_i} \big\| < \frac{1}{k} + \big\|  \sum_{i\in F_k}{x_i} \big\|$$
			
			Con $k$ suficientemente grande se tiene que $\big\| \displaystyle \sum_{i \in F_k}{x_i} \big\| > \tilde\varepsilon$ para cierto $\tilde \varepsilon >0$
			
			Tomando la permutación $\pi$ que deja fijos los elementos de $F_k$ se tiene que, para $k$ suficientemente grande 
			$$\varepsilon < \big\|  \sum_{i=1}^{n_k}{x_{\pi(i)}} + \sum_{i \in F_k}{x_i} - x \big\| \geq   \big\| \sum_{i \in F_k}{x_i} \big\| - \big\| \sum_{i=1}^{n_k}{x_{\pi(i)}}  - x\big\| > \tilde\varepsilon -  \big\|  \sum_{i=1}^{n_k}{x_{\pi(i)}}  - x\big\|$$
			Lo cual es una contradicción, pues suponiendo que la serie $\sum{x_{\pi(i)}}$ converge entonces el término $\big\| \displaystyle \sum_{i=1}^{n_k}{x_{\pi(i)}}  - x\big\|$ tiende a 0, pero entonces 
			$$\big\| \sum_{i=1}^{n_k}{x_{\pi(i)}} + \sum_{i \in F_k}{x_i} - x \big\| \not\longrightarrow 0$$
			contradiciendo la hipótesis de partida.
			
		\end{proof}
	
	
	
	%%%%%%%%%%%%%%%%%%%%%%%%%%%%%%%%%%%%%%%%%%%%%%%%%%%%%%%%%%%
	
	%Por simplicidad voy a fusionar esta definición con el lema siguiente para ahorrar espacio
	
	% \begin{definition}
		% Se dice que $\sum{x_i}$ es una serie incondicionalmente de Cauchy si dado $\varepsilon > 0$ existe un subconjunto finito $F$ de $\mathbb{N}$ tal que 
		% $\big\| \displaystyle\sum_{i \in F'}{x_i} \big\|< \varepsilon$
		% para cualquier $F'$ conjunto finito disjunto de $F$.
		% \end{definition}
	
	%%%%%%%%%%%%%%%%%%%%%%%%%%%%%%%%%%%%%%%%%%%%%%%%%%%%%%%%%%%
	
	Sea $(X,\|\cdot\|)$ un espacio de Banach.
	
	\begin{lemma}\label{lm:incon-cauchy}		
		Una serie $\sum{x_i}$ es incondicionalmente convergente si y solo si es incondicionalmente de Cauchy, i.e., si dado $\varepsilon > 0$ existe un subconjunto finito $F$ de $\mathbb{N}$ tal que 
		$$\big\| \sum_{i \in F'}{x_i} \big\|< \varepsilon$$
		para cualquier $F'$ conjunto finito disjunto de $F$.
		
	\end{lemma}	
		
		\begin{proof}
			Primero es necesario ver que una serie incondicionalmente de Cauchy es una serie de Cauchy.
			
			Fijado $\varepsilon >0$ y con $F$ un conjunto finito dado por la definición de incondicionalmente de Cauchy basta tomar $N = \max F$ con lo que tiene que para todo par $p,q > N$, 
			\[\big\| \sum_{i=1}^{p}{x_i} - \sum_{i=1}^{q}{x_i}  \big\| = \big\| \sum_{i=p}^{q}{x_i} \big\| < \varepsilon\]
			tomando como $F' =\{p,\ldots, q\}$, el cual es disjunto de $F$. Y por tanto, se cumple la definición de serie de Cauchy. 
			
			Como estamos en un espacio completo, por ser $X$ de Banach, se tiene que existe un cierto $x \in X$ tal que $x = \sum{x_i}$. 
			
			\noindent ``$\Longrightarrow$''
			
			Fijado un $\varepsilon > 0$, por ser condicionalmente de Cauchy se tiene que existe un $F^*$ subconjunto finito de $\mathbb{N}$ tal que para todo $F'$ subconjunto finito disjunto de $F^*$ se tiene que 
			\[\big\| \displaystyle \sum_{i \in F'}{x_i} \big\| < \varepsilon.\]
			Por la convergencia de la serie se tiene que existe un cierto $N_1$ natural para el cual, dado un $n>N_1$,
			\[\big\| \sum_{i=1}^{n}{x_i} - x \big\| < \varepsilon.\]
			
			Si se toma $n_0 = \max \{ \max F^*, N_1\}$ y $F= \{1, \ldots, n_0\}$ se tiene que, para cualquier $F' \supset F$  
			\[\big\| \sum_{i \in F'}{x_i} - x \big\| = \big\| \sum_{i=1}^{n_0}{x_i} + \sum_{i \in F', i > n_0}{x_i} - x\big\| \leq \big\| \sum_{i=1}^{n_0}{x_i} -x\big\| + \big\| \sum_{i \in F', i > n_0}{x_i}\big\| < \varepsilon + \varepsilon = 2\varepsilon\]
			
			Y por tanto se tiene la convergencia incondicional.
			
			\noindent ``$\Longleftarrow$''
			
			Fijado $\varepsilon >0$, por la convergencia incondicional se tiene que existe $F$ conjunto finito de $\mathbb{N}$ de forma que dado un conjunto finito $F' \supset F$ se tiene que 
			$$\big\| \sum_{i \in F'}{x_i} - x \big\| < \varepsilon$$
			
			de la misma forma, por la convergencia de la serie se tiene que existe un cierto $N_1$ natural para el cual, dado un $n>N_1$ 
			\[\big\| \sum_{i=1}^{n}{x_i} - x \big\| < \varepsilon.\]
			Si se toma $n_0 = \max \{ \max F, N_1\}$, dado un conjunto finito $F'$ disjunto de $F$ se tiene que $F \subset F' \cup \{1, \ldots, n_0\}$ y por tanto
			\[\big\| \sum_{i\in F'}{x_i} \big\| = \big\| \sum_{i \in F'}{x_i} + \sum_{i=1}^{n_0}{x_i} - x + x - \sum_{i=1}^{n_0}{x_i} \big\| \leq \big\| \sum_{i \in F' \cup \{1,\ldots,n_0\}}{x_i}  - x \big\| + \big\| x - \sum_{i=1}^{n_0}{x_i} \big\| < \varepsilon + \varepsilon = 2\varepsilon\]
			
		\end{proof}
	
	\begin{lemma} \label{lm:incondicional-epsilon}
		Una serie $\sum{x_i}$ converge incondicionalmente si y solo si converge $\sum{\varepsilon_i x_i}$ con cualquier elección de $(\varepsilon_i )_i= (\pm1)_i$.
	\end{lemma}
	
		\begin{proof}
			\noindent ``$\Longrightarrow$''
			
			% Incondicionalmente convergente implica condicionalmente cauchy que implica probar que la serie $\sum{\varepsilon_i x_i}$ es de cauchy empleando la idea que mostré en la nota anterior a la anterior.
			
			Probemos que $\sum \varepsilon_ix_i$ es de Cauchy y, por ser $X$ un espacio completo, convergente.
			
			Por ser $\sum x_i$ incondicionalmente convergente, por el Lema \ref{lm:incon-cauchy}, fijado $\varepsilon > 0$, existe un conjunto $F$ finito de $\mathbb{N}$ tal que para cualquier $F'$ disjunto de F que verifica 
			\[ \big\|  \sum_{i \in F'}{x_{i}}  \big\| < \varepsilon \]
			Sea $n_{0} > \max F$, dados $n,m > n_{0}$, se tiene que $\{i \in \mathbb{N} : \varepsilon_{i} = 1, m> i > n \} \cap F = \emptyset$ y $\{ i \in \mathbb{N} : \varepsilon_{i} = -1, m>i>n \}  \cap F = \emptyset$, por tanto:
			$$ \big\|  \sum_{i=n}^{m}{\varepsilon_{i}x_{i}}  \big\| \leq \big\|  \sum_{n \leq i \leq m, \varepsilon_{i}=1}{x_{i}}  \big\|\big\| \sum_{n \leq i \leq m, \varepsilon_{i}=-1}{- x_{i}}  \big\| < \varepsilon + \varepsilon = 2\varepsilon $$
			y por tanto la serie converge.
			
			
			\noindent ``$\Longleftarrow$''
			
			Suponiendo que $\sum{x_{i}}$ no es incondicionalmente de Cauchy, entonces existe $\varepsilon>0$ tal que para cualquier subconjunto de $\mathbb{N}$ finito $F$ existe otro conjunto finito $F'$ disjunto de $F$ que verifica 
			\[\big\|  \sum_{i \in F'}{x_{i}}  \big\| > \varepsilon\]
			En particular, si se toma $F^{*} = \{ 1 \}$ finito, existe $F_{1}$ disjunto tal que $\big\| \displaystyle \sum_{i \in F_{1}}{x_{i}}  \big\| > \varepsilon$.
			Si ahora se toma $F^{*} = \{ 1, \dots, \max F_{1} \}$ finito se tiene que existe un $F_{2}$ disjunto de $F^*$ verificando que 
			\[\big\|  \displaystyle \sum_{i \in F_{2}}{x_{i}}  \big\| > \varepsilon\]
			y que $\min F_{2} > \max F_{1}$.
			
			Por recurrencia se definen todos los $F_{k}$ con $k$ natural de forma que $\big\|  \displaystyle \sum_{i \in F_{k}}{x_{i}}  \big\| > \varepsilon$ y que $\min F_{k+1} > \max F_{k}$.
			
			Se definen ahora los conjuntos $E_{k} = \{ \min F_{k}, \dots, \max F_{k} \}$ para todo $k$ natural y los valores $v_i$ como
			\begin{itemize}
				\item[$\circ$] $v_{i} = 1$ si $i \in \bigcup F_{k}$
				\item[$\circ$] $v_{i}= -1$ si $i \not\in \bigcup F_{k}$.
			\end{itemize}
			De esta forma,
			%%%%%%%
			\[\big\|  \sum_{i \in E_{k}}{x_{i}}  \big\| + \big\|  \sum_{i \in E_{k}}{v_{i}x_{i}}  \big\| \geq \big\|  \sum_{i \in E_{k}}{x_{i}} + \sum_{i \in E_{k}}{v_{i}x_{i}}  \big\| = \big\|  \sum_{i \in F_{k}}{x_{i}}  \big\| > 2 \varepsilon\]
			%%%%%%
			Por tanto, o bien $\big\| \displaystyle \sum_{i \in E_{k}}{x_{i}}  \big\| > \varepsilon$, y en ese caso se tomará $\varepsilon_{i}=1$ para todo $i \in E_{k}$, o bien $\big\| \displaystyle \sum_{i \in E_{k}}{v_{i}x_{i}}  \big\| > \varepsilon$, y en ese caso $\varepsilon_{i} = v_{i}$ para todo $i \in E_{k}$. En cualquier caso, se asegura que 
			%%%%%%
			\[\big\|  \sum_{i \in E_{k}}{\varepsilon_{i}x_{i}}  \big\| \geq \varepsilon\]
			y por tanto la serie no es de Cauchy, pues para todo $N$ natural existe un $k$ natural de forma que $p = \min F_{k}, q = \max F_{k} > N$:
			\[\big\|  \sum_{i=p}^{q}{\varepsilon_{i}x_{i}}  \big\| = \big\|  \sum_{i \in E_{k}}{\varepsilon_{i}x_{i}}  \big\| > \varepsilon.\]
			
		\end{proof}
	
	\begin{lemma}\label{lm:inc-conv-conv-subset}
		Una serie $\sum{x_{i}}$ es incondicionalmente convergente si y solo si $\sum{x_{n_{i}}}$ es convergente para cualquier sucesión creciente $\{ n_{i} \}_{i=1}^{\infty}$.
		
	\end{lemma}
	
		\begin{proof}
			
			\noindent ``$\Longrightarrow$''   
			
			Por ser incondicionalmente convergente (y por tanto incondicionalmente de Cauchy) existe $F$ conjunto finito de $\mathbb{N}$ tal que para cualquier $F'$ conjunto finito disjunto se tiene que 
			\[\big\| \displaystyle \sum_{i \in F'}{x_{i}}  \big\| < \varepsilon.\]
			En particular, si $p,q > \max F$, entonces $n_{p}, n_{q} > \max F$ y 
			\[\big\|  \sum_{i \in =p}^{q}{x_{n_{i}}}  \big\| < \varepsilon.\]
			Por tanto la serie es de Cauchy, ergo convergente.
			
			\noindent ``$\Longleftarrow$''
			
			Viendo el contrarrecíproco, supongamos que la serie no es incondicionalmente de Cauchy, entonces existe un $\varepsilon >0$ y una sucesión de conjuntos finitos $F_{k}$ de $\mathbb{N}$ tal que $\max F_{k} < \min F_{k+1}$ y 
			\[ \big\| \sum_{i \in F_{k}}{x_{i}}  \big\| \geq \varepsilon.\]
			Tomando por $\{ n_{i} \}_{i=1}^{\infty} = \bigcup F_{k}$, se tiene que, para  cualquier $N$ natural existe un $k$ tal que $\min F_{k}, \max F_{k} > N$:
			\[\big\| \sum_{i = \min F_{k}}^{\max F_{k}}{x_{n_{i}}}  \big\| \geq \varepsilon\]
			y por tanto la serie no es de Cauchy. \gls{latex}
			
		\end{proof}
	
	
	
	%%%%%%%%%%%%%%%%%%%%%%%%%%%%%%%%%%%%%%%%%%%%%%%%%%%%%%%%%%%%%%%%%%%%%%%%%%%%
	% \begin{lemma}
		
		% Dada una serie $\sum{x_{i}}$ incondicionalmente convergente, para todo $\varepsilon > 0$ existe un $i_{0}$ tal que 
		% $$\big\| \displaystyle \sum_{i = i_{0}}^{\infty}{\varepsilon_{i}x_{i}}  \big\| < \varepsilon$$
		% para cualquier elección de $\varepsilon_{i} = \pm 1$.
		
		% \begin{proof}\hspace{1cm}
			
			% Por el Lema \ref{lm:incondicional-epsilon}, como la serie es incondicionalmente convergente entonces $\sum{\varepsilon_{i}x_{i}}$ es convergente para cualquier elección de $\varepsilon_{i}= \pm 1$, y equivalentemente de Cauchy. Por tanto, existe un cierto $i_{0}$ tal que para cualquier $j>i_{0}$ se tendrá
			% $$\big\|  \sum_{i =i_{0}}^{j}{x_{i}}  \big\| < \varepsilon.$$
			% Tomando límites en $j \rightarrow \infty$, como la norma es continua
			% $$\big\| \sum_{i =i_{0}}^{\infty}{x_{i}}  \big\| < \varepsilon.$$
			
			% \end{proof}
		%\end{lemma}
		
		% \begin{lemma}
			
			% Dada una serie $\sum{x_{i}}$ incondicionalmente convergente, entonces el conjunto de vectores de la forma $\sum{\varepsilon_{i}x_{i}}$ con $\varepsilon_{i} = \pm 1$ es un compacto de $X$.
			
			
			% \begin{proof}\hspace{1cm}
				
				% Sea $\pi : \{ 1,-1 \}^{\mathbb{N}} \longrightarrow X$ la aplicación que a cada selección de $\{ \varepsilon_{i} \}$ le asigna $\sum{\varepsilon_{i}x_{i}}$. Si se demuestra que $\pi$ es continua y que $\{ 1,-1 \}^{\mathbb{N}}$ es un compacto con la topología de la convergencia puntual.
				
				% Véase que $\pi$ es continua. Por la incondicionalidad de Cauchy se tiene que para todo $\varepsilon >0$ existe un $n_{0}$ suficientemente grande tal que para cualquier $F'$ conjunto finito con $\min F' > n_{0}$ se verifica
				% $$\big\|  \sum_{i \in F'}{x_{i}}  \big\| < \varepsilon$$
				
				% Dadas dos $\{ \varepsilon_{i} \}$ y $\{ \tilde\varepsilon_{i} \}$ tales que $\varepsilon_{i} = \tilde\varepsilon_{i}$ para todo $i<n_{0}$ se tiene que:
				% $$\| \pi(\{ \varepsilon_{i} \}) - \pi(\{ \tilde \varepsilon_{i} \}) \| = \big\|  \sum{\varepsilon_{i}x_{i}} - \sum{\tilde\varepsilon_{i}x_{i}}  \big\| = \lim_{  j \to \infty }{\big\|  \sum_{i=1}^{j}{\varepsilon_{i}x_{i}} - \sum_{i=1}^{j}{\tilde\varepsilon_{i}x_{i}}  \big\|} =  \lim_{  j \to \infty }{2\big\|  \sum_{i=n_{0}}^{j}{\delta_{i}x_{i}}  \big\|}$$
				% donde $\delta_{i} = \varepsilon_{i}$,  si  $\varepsilon_{i}\tilde \varepsilon_{i}=1$ y $\delta_{i} = 0$ en otro caso. De esta forma,
				% $$\| \pi(\{ \varepsilon_{i} \}) - \pi(\{ \tilde \varepsilon_{i} \}) \| = \lim_{  j \to \infty }{2\big\|  \sum_{i=n_{0}}^{j}{\delta_{i}x_{i}}  \big\|} \leq \lim_{ j \to \infty }{2 \big [ \big\|  \sum_{\varepsilon_{i} = 1}{x_{i}}  \big\| + \big\|  \sum_{\varepsilon_{i} =-1}{x_{i}}  \big\| \big]} \leq 4 \varepsilon$$
				% pues $\{ i \in \mathbb{N} : \varepsilon_{i}=1, j>i >n_{0}\}$ y $\{ i \in \mathbb{N} : \varepsilon_{i}=-1, j>i >n_{0}\}$ son conjuntos finitos en las hipótesis de la incondicionalidad de Cauchy. Por tanto, la aplicación $\pi$ es continua.  
				
				% \end{proof}
			% \end{lemma}
		
		
		% \begin{theorem}
			
			% Dada una serie $\sum{x_{i}}$ incondicionalmente convergente,y el conjunto 
			% $$S = \big\{  \sum_{i=1}^{n}{x_{i}} : n \in \mathbb{N}, \varepsilon_{i}= \pm 1  \big\},$$ S es relativamente compacto en $X$.
			
			
			% \begin{proof}
				
				% Denotando al conjunto de los vectores de la forma $\sum{\varepsilon_{i}x_{i}}$ con $\varepsilon_{i} = \pm 1$ por $F$, basta ver que $\bar{S} \subset \left( \frac{1}{2}F+\frac{1}{2}F \right) \cup F$. Como $F$ es compacto y el producto por escalares es continuo, el conjunto $\frac{1}{2}F+\frac{1}{2}F := \left\{  \frac{a}{2}+\frac{b}{2} : a \in F, b \in F  \right\}$ es compacto y la unión de compactos es compacto, y como $\bar{S}$ es un cerrado dentro de un compacto, sería también compacto. 
				
				% Véase ahora que, en efecto, $\bar{S} \subset \left( \frac{1}{2}F+\frac{1}{2}F \right) \cup F$. $\bar{S} = F \cup S$, pues basta tomar límite $n \rightarrow \infty$ para llegar a cualquier elemento de $F$ (basta considerar las sumas parciales). Sólo es preciso ver que $S \subset \frac{1}{2}F+\frac{1}{2}F$. El resultado se sigue de ver que dado un elemento $\displaystyle \sum_{i=1}^{n}{\varepsilon_{i}x_{i}}$ de S, se puede escribir como
				% $$
				%  \sum_{i=1}^{n}{\varepsilon_{i}x_{i}} = \frac{1}{2}\sum_{i=1}^{\infty}{\varepsilon_{i}x_{i}} + \frac{1}{2}\sum_{i=1}^{\infty}{\tilde\varepsilon_{i}x_{i}}
				% $$
				% donde $\varepsilon_{i} = \tilde\varepsilon_{i}$ para todo $i \in \{ 1,\dots,n \}$ y $\varepsilon_{i} = - \tilde\varepsilon_{i}$ para todo $i > n$.
				
				% \end{proof}
			
			%\end{theorem}
			%%%%%%%%%%%%%%%%%%%%%%%%%%%%%%%%%%%%%%%%%%%%%%%%%%%%%%%%%%%%%%%%%%%%%
			
			% \begin{definition}
				%     Dado un espacio de Banach $X$, definimos el conjunto 
				%     $$(\sum \oplus\, X)_{\ell_p}$$
				%     con $p\in [1,\infty)$ como el conjunto de sucesiones $(x_i)_i$ con términos en $X$ tales que 
				%     $$(\sum\|x_i\|^p)^\frac{1}{p} <  \infty$$
				%     es decir, una generalización del espacio $\ell_p$ a conjuntos distintos de los reales.
				
				%     De la misma forma, definimos el conjunto
				%     $$(\sum \oplus\, X)_{c_0}$$
				%     como el conjunto de sucesiones $(x_i)_i$ tales que 
				%     $$\sup_{i \in \N}\|x_i\| < \infty$$
				% \end{definition}
			
			% En particular, cuando tenemos que $X$ es $\ell_p$ o $c_0$ tenemos el siguiente resultado:
			
			% \begin{theorem}
				%     Denotando por $X$ al espacio $c_0$ o $\ell_p$ con $p\in[1,\infty)$, tenemos que 
				%     $$\left(\sum \oplus\, X\right)_{X} \simeq X$$
				
				%     \begin{proof}
					
					%         La idea es la misma que la prueba de $\N\times \N$ es biyectivo con $\N$.
					
					%         La cuestión se reduce realmente a encontrar una biyección $f\colon \N \times \N \to \N$ de forma que el isomorfismo $F\colon \left(\sum \oplus\, X\right)_{X} \to X$ se defina como $F(((x_i^n)_n)_i) = (y_{m})_m$ donde $m = f(i,n)$ y $y_m = x_i^n$. Obviamente $F$ es biyectiva pues si dos elementos $x$ y $z$ de $\left(\sum \oplus\, X\right)_{X}$ tienen la misma imagen entonces son iguales para todas sus coordenadas:
					%         $$z_i^n = y_m = y_{f(i,n)} = x_i^n \quad \forall i,n \in \N$$
					%         y para toda sucesión $(x_m)_m$ de $X$ podemos ver su preimagen como $y^i_n = x_m$ con $(i,n) = f^{-1}(m)$.
					
					%         La linealidad resulta fácil de probar, de nuevo tomando dos elementos $x,z \in \left(\sum \oplus\, X\right)_{X}$ y $\lambda\in\R$ tenemos que 
					%         $$F(\lambda x + z) = F(((\lambda x_i^n+z_i^n)_n)_i) = (y_m)_m \quad \text{donde }\, y_{f(i,n)} = \lambda x_i^n + z_i^n   $$
					%         $$(y_m)_m = \lambda(y^1_m)_m+ (y^2_m)_m = \lambda F(x) + F(y) \quad \text{donde }\, y^1_{f(i,n)} = x_i^n;\, y²_{f(i,n)} = z_i^n$$
					
					%         Para la continuidad nos bastará ver que la norma de $F(x)$ coincide con la de $x$. 
					
					%         Para el caso de $\ell_p$:
					
					%         Dado un elemento $x = (x_i)_i = ((x_i^n)_n)_i \in\left(\sum \oplus\, \ell_p\right)_{\ell_p}$, sabemos que 
					%         \begin{align*}
						%         \|x\|_{\ell_p} &= \left(\sum_{i = 1}^\infty \|(x_i^n)_n\|^p_p\right)^\frac{1}{p} = \left(\sum_{i = 1}^\infty \sum_{n = 1}^\infty|x_i^n|^p\right)^\frac{1}{p} = \left(\sum_{i = 1}^\infty \sum_{n = 1}^\infty|y_{f(i,n)}|^p\right)^\frac{1}{p} \\&= \left(\sum_{m = 1}^\infty |y_m|^p\right)^\frac{1}{p} = \|F(x)\|_p
						%         \end{align*}
					%         y por tanto $F$ es continua.
					
					%         Para el caso de $c_0$:
					
					%         Dado un elemento $x = (x_i)_i = ((x_i^n)_n)_i \in\left(\sum \oplus\, c_o\right)_{c_0}$, sabemos que
					%         $$\|x\|_{c_0} = \sup_{i \in \N}\|(x_i^n)_n\|_\infty = \sup_{i \in \N}\sup_{n \in \N} |x_i^n| = \sup_{i \in \N}\sup_{n \in \N} |y_{f(i,n)}| =\sup_{m \in \N} |y_m|=  \|F(x)\|_\infty$$
					%         y por tanto $F$ es continua.
					%     \end{proof}
				% \end{theorem}
			
			% Como caso particular, si en vez de realizar una suma directa infinita realizamos una suma directa finita la demostración sería totalmente análoga pero empleando un isomorfismo basado en una biyección entre $\N$ y $\N\times\{0,\dots,n\}$ donde $n$ es la cantidad de sumandos.
			
			% Incluimos también el resultado para el caso de dos sumandos por emplearlo explícitamente en este trabajo.
			
			\begin{lemma}\label{lm:X-XX}
				Denotando por $X$ al espacio $c_0$ o $\ell_p$ con $p\in [1,\infty)$ tenemos que 
				\[X\oplus X \simeq X\]
			\end{lemma}
			
				\begin{proof}
					Primero hay que dotar al espacio suma de la norma producto:
					\[\|x+ y\| = \|x\|_X + \|y\|_X \quad \text{donde }\, x,y \in X\]
					
					Para probar el isomorfismo veamos que la aplicación $T\colon X\oplus X \to X$ definida como 
					\[T((x_n)_n,(y_n)_n)  = (x_1, y_1,x_2,y_2,x_3,y_3,\dots)\]
					es un isomorfismo. Por completitud, la expresión general sería
					\[T((x_n)_n,(y_n)_n) = (z_m)_m \quad \text{donde }\, z_m = \left\{\begin{aligned} x_n, &\quad m = 2n \\ y_n, &\quad m = 2n+1\end{aligned}\right.\]
					
					La linealidad es fácil de probar: 
					
					dados dos elementos $((x_n)_n,(y_n)_n), ((a_n)_n,(b_n)_n) \in X\oplus X$ y $\lambda \in \R$:
					\[T(\lambda((x_n)_n,(y_n)_n) +  ((a_n)_n,(b_n)_n)) = T(((\lambda x_n+ a_n)_n,(\lambda y_n + b_n)_n)) = (z_m)_m\]
					donde, para todo $m$ natural
					\[z_m = \left\{\begin{aligned} \lambda x_n + a_n, &\quad m = 2n \\ \lambda y_n + b_n, &\quad m = 2n+1\end{aligned}\right. \]
					Por otro lado, 
					\[\lambda T(((x_n)_n,(y_n)_n))  = \lambda (z^1_m)_m \quad \text{donde }\, z_m^1 = \left\{\begin{aligned} x_n, &\quad m = 2n \\ y_n, &\quad m = 2n+1\end{aligned}\right. \]
					\[T(((a_n)_n,(b_n)_n)) = (z^2_m)_m \quad \text{donde }\, z^2_m = \left\{\begin{aligned} a_n, &\quad m = 2n \\ b_n, &\quad m = 2n+1\end{aligned}\right.\]
					Resulta por tanto evidente que 
					\[T(\lambda((x_n)_n,(y_n)_n) +  ((a_n)_n,(b_n)_n))  =\lambda T(((x_n)_n,(y_n)_n)) + T(((a_n)_n,(b_n)_n)) \]
					
					La biyectividad es sencilla, Para toda sucesión de $X$ su preimagen es coger las subsucesiones de índices pares e índices impares, por tanto $T$  es sobreyectiva (dichas subsucesiones están en $X$ de forma trivial). La inyectividad es trivial.
					
					Por último, la continuidad.
					
					Dado un elemento $((x_n)_n,(y_n)_n) \in X\oplus X$,
					
					\begin{itemize}
						\item Para el caso de $X = c_0$,
						\begin{align*}
							\|T(((x_n)_n,(y_n)_n))\|_\infty & = \sup_{n \in \N}\{|x_n|,|y_n|\} \leq \sup_{n \in \N}|x_n| + \sup_{n \in \N}|y_n| \\&= \|(x_n)_n\|_\infty+\|(y_n)_n\|_\infty = \|((x_n)_n,(y_n)_n)\|
						\end{align*}
						\item Para el caso de $\ell_p$
						\begin{align*}
							\|T(((x_n)_n,(y_n)_n))\|_p & = \big(\sum_{n = 1}^\infty |x_n|^p + |y_n|^p\big)^\frac{1}{p} = \big(\sum_{n = 1}^\infty |x_n|^p +\sum_{n =1}^\infty |y_n|^p\big)^\frac{1}{p}\\ &\leq \big(\sum_{n = 1}^\infty |x_n|^p\big)^\frac{1}{p} + \big(\sum_{n = 1}^\infty|y_n|^p\big)^\frac{1}{p} = \|(x_n)_n\|_p+\|(y_n)_n\|_p \\&= \|((x_n)_n,(y_n)_n)\|
						\end{align*}
					\end{itemize}
					y por tanto se da que $T$ es continua, demostrándose que $T$ es un isomorfismo. 
				\end{proof}
			
			
			\newpage
			\section*{Teoremas fundamentales del Análisis funcional}
			
			% para poner los teoremas relevantes pero que no son materia del TFG per se.
			
			
			
			\begin{theorem}[de la aplicación abierta{ \cite[pág. 35]{Brezis}}]\label{th:ap-abierta}
				Dados $X,Y$ dos espacios de Banach y $f:X \rightarrow Y$ un operador lineal continuo y sobreyectivo, entonces $f$ es una aplicación abierta.
			\end{theorem}
			% debería incluír alguna referencia a donde ver este teorema. Tal vez el rudin?
			
			\begin{theorem}[del grafo cerrado {\cite[pág. 37]{Brezis}}]\label{th:grafo-cerrado}
				Sean $X,Y$ dos espacios de Banach y $T\colon X \to Y$ un operador lineal. Si para toda sucesión $(x_n)_n$ de $X$ tal que $x_n\xrightarrow{n\to\infty}x$ se verifica que, si $T(x_n) \xrightarrow{n\to \infty}y$ entonces $y = T(x)$, entonces $T$ es un operador continuo.
			\end{theorem}
			
			\begin{theorem}[de extensión de Hahn-Banach]\label{th:ext-hahn-Banach}
				Dados $X$ espacio normado, $Y$ espacio de Banach y $M$ un subespacio denso en $X$. Si $T_0 \colon M \to Y$ es un operador lineal continuo entonces existe un único $T:X \to Y$ operador lineal acotado que verifica
				\begin{itemize}
					\item[(i)] $T|_{M} = T_0$
					\item[(ii)] $\|T\| = \|T_0\|$ 
				\end{itemize}
			\end{theorem}
			
			\begin{theorem}[de extensión de Hahn-Banach aplicado a funcionales]\label{th:Hahn-Banach-funcionales}
				Sea $X$ un espacio normado e $Y$ un subespacio de $X$. Si $f_0\in Y^*$, entonces existe un $f\in X^*$ tal que 
				\begin{itemize}
					\item[$(i)$] $\|f_0\| = \|f\|$
					\item[$(ii)$] $f|_Y = f_0$ 
				\end{itemize}
				
			\end{theorem}
			
			\begin{lemma}[de Riesz de elementos casi ortogonales {\cite[pág. 2]{SequencesSeriesDiestel}}]\label{lm:riesz}
				Dado un espacio normado $X$ y un subespacio $Y$ no denso en $X$ entonces para todo $r \in (0,1)$ existe un cierto $x \in X$ con $\|x\| = 1$ de forma que
				$$d(x,Y) \geq r$$
				
			\end{lemma}
			
			\begin{theorem}[Principio de acotación uniforme de Banach-Steinhaus {\cite[pág. 32]{Brezis}}]\label{th:acotacion-uniforme-banach-steinhaus}
				Dados $X,Y$ espacios normados y $\mathcal{F}$ una familia de operadores lineales de $X$ a $Y$. Si $\sup_{T\in\mathcal{F}}\|Tx\|$ para todo $x\in X$, entonces $\sup_{T\in \mathcal{F}}\|T\|<\infty$.
			\end{theorem}
	
	
%	\chapter{Prueba de los anexos}
%	
%	\begin{definicion}{Moda Variable por el tiempo invertido en la toma de contacto de los dos años luz que tardamos en completarlo}
%		Dada una variable aleatoria $Y$, absolutamente continua, se definen
%		sus modas como aquellos valores $y\in\operatorname{Sop}(Y)$ que son
%		máximos locales de su función de densidad.
%		\begin{equation}
%			\label{eq:2}
%			1+1 = 2
%		\end{equation}
%		Vamos a ver si las referencias seguirían funcionando de momento: ecuación \ref{eq:2}
%	\end{definicion}
%	
%	Fíjate en que la numeración de definiciones/teoremas sigue siendo
%	plana y continua (no se reinicia en el anexo ni depende de que
%	\thechapter{} ahora sea una letra).
%	

\begin{thebibliography}{2}\addcontentsline{toc}{chapter}{Bibliografía}
	
	\bibitem{Brezis}
	Leslie Lamport,
	\textit{\LaTeX: a document preparation system},
	Addison-Wesley, Reading, Massachusetts,
	2nd edition, 1994. 
	
	
	
	\bibitem{einstein05}
	Albert Einstein,
	\textit{Zur Elektrodynamik bewegter K{\"o}rper},
	Annalen der Physik, 17:891--921, 1905.
	
\end{thebibliography}




% Para modificar el texto del glosario
\renewcommand{\glossarypreamble}{\par En la siguiente sección se presentan diferentes conceptos y acrónimos, junto con sus definiciones, que son necesarios para el entendimiento completo del texto.}

\printglossaries





\end{document}